\documentclass[11pt]{article}
\usepackage{style}

%COMMENTS
\newcommand{\Anouksays}[1]{\textcolor{orange}{\textbf{\underline{Anouk says:}} #1}}
\newcommand{\Gyivansays}[1]{\textcolor{blue}{\textbf{\underline{Gyivan says:}} #1}}
\newcommand{\Tonosays}[1]{\textcolor{purple}{\textbf{\underline{To\~no says:}} #1}}
\newcommand{\jdlsays}[1]{\textcolor{lightgreen}{\textbf{\underline{Jesús says:}} #1}}


\newcommand{\ba}{\mathbf{a}}
\newcommand{\bb}{\mathbf{b}}
\newcommand{\bc}{\mathbf{c}}
\newcommand{\be}{\mathbf{e}}
\newcommand{\bp}{\mathbf{p}}
\newcommand{\bs}{\mathbf{s}}
\newcommand{\bt}{\mathbf{t}}
\newcommand{\bu}{\mathbf{u}}
\newcommand{\bv}{\mathbf{v}}
\newcommand{\bw}{\mathbf{w}}
\newcommand{\bx}{\mathbf{x}}
\newcommand{\by}{\mathbf{y}}
\newcommand{\bz}{\mathbf{z}}


\newcommand{\R}{\mathbb{R}}
\newcommand{\Z}{\mathbb{Z}}
\newcommand{\N}{\mathbb{N}}
\newcommand{\Q}{\mathbb{Q}}
\newcommand{\C}{\mathbb{C}}
\renewcommand{\P}{\mathbb{P}}

\newcommand{\Lcal}{\mathcal{L}}
\newcommand{\Dcal}{\mathcal{D}}
\newcommand{\Ccal}{\mathcal{C}}
\newcommand{\Scal}{\mathcal{S}}
\newcommand{\Tcal}{\mathcal{T}}
\newcommand{\Rcal}{\mathcal{R}}

\DeclareMathOperator{\argmin}{argmin}

\DeclareMathOperator{\intr}{int}
\DeclareMathOperator{\relintr}{relint}
\DeclareMathOperator{\conv}{conv}
\DeclareMathOperator{\aff}{aff}
\DeclareMathOperator{\vol}{vol}
\DeclareMathOperator{\spn}{span}
\DeclareMathOperator{\dir}{dir}
\DeclareMathOperator{\ldiam}{diam_\Z}
\DeclareMathOperator{\diam}{diam}
\DeclareMathOperator{\ldist}{dist_\Z}
\DeclareMathOperator{\nvol}{nvol}
\DeclareMathOperator{\verts}{vert}
\DeclareMathOperator{\edges}{edg}

\DeclareMathOperator{\blocks}{blocks}
\DeclareMathOperator{\rem}{rem}
\DeclareMathOperator{\LD}{LD}
\DeclareMathOperator{\Ehr}{L}

% %WITH PARENTHESES
% \newcommand{\intr}[1]{\operatorname{int}\left(#1\right)}
% \newcommand{\relintr}[1]{\operatorname{relint}\left(#1\right)}
% \newcommand{\conv}[1]{\operatorname{conv}\left(#1\right)}
% \newcommand{\aff}[1]{\operatorname{aff}\left(#1\right)}
% \newcommand{\vol}[1]{\operatorname{vol}\left(#1\right)}
% \newcommand{\spn}[1]{\operatorname{span}\left(#1\right)}
% \newcommand{\dir}[1]{\operatorname{dir}\left(#1\right)}
% \newcommand{\ldiam}[1]{\operatorname{diam}_{\mathbb{Z}}\left(#1\right)}
% \newcommand{\diam}[1]{\operatorname{diam}\left(#1\right)}
% \newcommand{\ldist}[1]{\operatorname{dist}_{\mathbb{Z}}\left(#1\right)}
% \newcommand{\nvol}[1]{\operatorname{nvol}\left(#1\right)}
% \newcommand{\verts}[1]{\operatorname{vert}\left(#1\right)}
% \newcommand{\edges}[1]{\operatorname{edg}\left(#1\right)}

\DeclareMathOperator{\GL}{GL} %general linear group
\DeclareMathOperator{\BG}{\mathsf{BG}} %Borsuk Graph Gyivan font
\DeclareMathOperator{\NP}{\mathsf{NP}} %Borsuk Graph Gyivan font

%\DeclareMathOperator{\blocks}{blocks} 
%\DeclareMathOperator{\rem}{rem}

\newcommand{\sgn}{sign} %sign
\newcommand{\st}{\, \mathrm{:}\,}


\newcommand{\lp}{\left(}
\newcommand{\rp}{\right)}
\newcommand{\lb}{\left\{}
\newcommand{\rb}{\right\}}

\newcommand{\lvec}[1]{\overleftarrow{#1}}

%COLORS ANOUK USES: darklav, lavender, orange (normal orange), lightgreen on first page
%Or lavender!90!gray lavender!50!black or white etc. for 90% lavender, or 50% lavender

%USED
\definecolor{darklav}{rgb}{0.45, 0.31, 0.59}
\definecolor{lightgreen}{rgb}{0.53, 0.66, 0.42} %``asparagus''
\definecolor{etonblue}{rgb}{0.59, 0.78, 0.64} %light turqoise-ish
\definecolor{frenchblue}{rgb}{0.0, 0.45, 0.73} %my favorite blue
\definecolor{lavender}{rgb}{0.84, 0.79, 0.87}
\definecolor{lilac}{rgb}{0.78, 0.64, 0.78} %light pale purple
\definecolor{darkyellow}{rgb}{1.0, 0.75, 0.0} %``amber''
\definecolor{redgrapefruit}{rgb}{1.0, 0.44, 0.37} %``bittersweet''
\definecolor{brightube}{rgb}{0.82, 0.62, 0.91} %light pale & bright purple
\definecolor{stronggreen}{rgb}{0.0, 0.42, 0.24} %``cadium green''

%NOT USED
\definecolor{darkpurp}{rgb}{0.36, 0.22, 0.33}
\definecolor{eggplant}{rgb}{0.38, 0.25, 0.32}
\definecolor{frenchlilac}{rgb}{0.53, 0.38, 0.56}
\definecolor{mauvetaupe}{rgb}{0.57, 0.37, 0.43}
\definecolor{stoneblue}{rgb}{0.45, 0.66, 0.76}
\definecolor{linen}{rgb}{0.98, 0.94, 0.9}

\title{\scshape On Lattice Diameter Segments and \mbox{A Discrete Borsuk Partition Problem}}
\date{}

% \author{Anouk E.~Brose\thanks{Corresponding author}, Jes\'us A.~De Loera, Gyivan Lopez-Campos, and Antonio J.~Torres}
        
\author[1]{Anouk E.~Brose \thanks{Corresponding author}} 
\author[1]{Jes\'us A.~De Loera}
\author[1]{Gyivan Lopez-Campos}
\author[1]{Antonio J.~Torres}

\affil[1]{\small \scshape Department of Mathematics, University of California, Davis}

% \address{
%  Department of Mathematics, University of California, Davis
% }
% \email{\{aebrose,XXX\}@ucdavis.edu}

\begin{document}

\maketitle

\mscclasses{\emph{Mathematics Subject Classification.}
52A38, %Length, area, volume and convex sets (aspects of convex geometry)
52B20, %Lattice polytopes (primary)
52B55, %Computational aspects related to convexity
52C07, %Lattice and convex bodes in R^n
52C17, %Packing and covering in n dimensions (for Borsuk)
68W35. %Computational complexity of algorithms
}

% 52C05, %Lattice and convex bodes in R^2
% 68R05, %Combinatorics in computer science
% 68W35 %Computational complexity of algorithms
% 52B55 %Computational aspects related to convexity

\keywords{\emph{Key words and phrases.} lattice diameter; lattice lines; lattice segments; lattice polygons; Ehrhart theory; computational complexity; Borsuk's partition problem.}

\begin{abstract}
    \noindent The lattice diameter of a bounded set $S \subset \R^d$ measures the maximal number of lattice points in a segment whose endpoints are lattice points in $S$. Such a segment is called a lattice diameter segment of $S$. This simple invariant yields interesting applications and challenges. We describe a polynomial-time algorithm that computes lattice diameter segments of lattice polygons and show that computing lattice diameters of semi-algebraic sets in dimensions three and higher is $\NP$-hard.
    We prove that the function that counts lattice diameter segments in dilations of a lattice polygon is eventually a quasi-polynomial in the dilation factor.
    % in dimension two we show that the count of lattice diameter segments under dilation agrees with a quasi-polynomial.
    We also study the number of \mbox{directions that} lattice diameter segments can have. Finally, we prove a Borsuk-type theorem on the number of parts needed to partition a set of lattice points such that each part has strictly smaller \mbox{lattice diameter}.
\end{abstract}

\section{Introduction}  

We study a fundamental geometric invariant of compact sets, called the \textit{lattice diameter}, with a variety of applications in discrete geometry, integer optimization, and the geometry of numbers. We are not the first to study lattice diameters.
This notion was introduced by Corzatt and \mbox{Stolarsky~\cite{Corzatt_PhD}, 
% who were interested in minimal covers of lattice polygons, i.e., the smallest number of lines that cover all lattice points in the polygon. 
%Corzatt's conjecture, which is still open, states that any lattice polygon has a minimal cover using lines in at most four different directions.
studied} by Alarcon~\cite{Alarcon_PhD} and later by Bárány and Füredi~\cite{Barany_Furedi}. Here are the key definitions of this paper: 

\begin{definition}\label{def:ld-non-convex}
    For a compact set $S \subset \R^d$, define $$\ldiam(S):=\max_{\bx,\by \in S \cap \Z^d} |\conv(\{\bx,\by\}) \cap \Z^d|-1$$
    as the \emph{lattice diameter} of $S$. If $\bx$ and $\by$ maximize $\ldiam(S)$, we call $\conv(\{\bx,\by\})$ a \emph{lattice diameter segment}. Similarly, $\aff(\{\bx,\by\})$ is a \emph{lattice diameter line}, and finally, $\by-\bx$ or any of its non-zero multiples is called a \emph{lattice diameter direction}. 
\end{definition}

Our paper investigates properties of lattice diameter segments and algorithms to compute them.

Note that the lattice diameter of $S$ and $\conv(S)$ can differ when $S$ is non-convex. When $S$ is a convex body, then a lattice diameter line is an affine lattice line $L$ such that $|L \cap S \cap \Z^d|$ is maximal among all affine lattice lines. In this case the lattice diameter of $S$ coincides with the lattice diameter of $\conv(\{S \cap \Z^d\})$, a lattice polytope, thus we focus on studying lattice diameters of lattice polytopes. The Euclidean diameters and lattice diameters of lattice polytopes need not correlate. For fixed lattice diameter, the Euclidean diameter can be arbitrarily large (\Cref{fig:unbounded_ed}). Moreover, Euclidean diameter lines always pass through vertices \cite[Lemma~1.7]{GritzmannKlee1992}, while lattice diameter lines \mbox{may not (\Cref{fig:LL_interior}).}
\begin{figure}[ht!]
  \centering

  \begin{subfigure}[t]{0.48\textwidth}
    \centering
    \begin{tikzpicture}[thick, scale=0.9, baseline=(current bounding box.north)]
      % unify bbox across BOTH subfigures
      \path[use as bounding box] (0,0) rectangle (7,5);

      \foreach \x in {0,...,7} {
        \foreach \y in {2,...,4} {
          \node at (\x,\y) [circle, draw=lightgray, fill=lightgray, minimum width=2pt] {};
        }
      }

      \draw[line width=0.3mm] (0,2) -- (5,3) -- (7,4) -- (2,3) -- cycle;
      \draw[line width=0.4mm, lightgreen] (0,2) -- (7,4);
      \draw[line width=0.3mm, orange]     (2,3) -- (5,3);

      \foreach \x in {2,...,5} \node at (\x,3) [circle, fill=orange, minimum width=4pt] {};

      \node at (1.1,3) [black]{$(0,0)$};
      \node at (5.9,3) [black]{$(\ell,0)$};
      \node at (0.5,1.6) [black]{$(-n,-1)$};
      \node at (6.5,4.4) [black]{$(\ell+n,1)$};
    \end{tikzpicture}
    \caption{Lattice polygons with fixed lattice diameter \mbox{(orange)} and \mbox{unbounded} Euclidean diameter (green).}
    \label{fig:unbounded_ed}
  \end{subfigure}\hfill
  \begin{subfigure}[t]{0.48\textwidth}
    \centering
    \begin{tikzpicture}[thick, scale=0.9, baseline=(current bounding box.north)]
      % unify bbox across BOTH subfigures
      \path[use as bounding box] (0,0) rectangle (7,5);

      \foreach \x in {0,...,7} {
        \foreach \y in {1,...,5} {
          \node at (\x,\y) [circle, draw=lightgray, fill=lightgray, minimum width=2pt] {};
        }
      }
      \draw[line width=0.3mm] (5,1) -- (7,4) -- (2,5) -- (0,2) -- (3,1) -- cycle;
      \draw[line width=0.3mm, orange] (1,3) -- (6,3);
      \foreach \x in {1,...,6} \node at (\x,3) [circle, fill=orange, minimum width=4pt] {};
    \end{tikzpicture}
    \caption{A lattice diameter segment contained in the interior of a lattice polygon.}
    \label{fig:LL_interior}
  \end{subfigure}

  \caption{Examples of lattice diameter segments in lattice polygons.}
  \label{fig:two_subfigs}
\end{figure}


% \begin{figure}[ht!]
%     \centering
%     % Left figure
%     \begin{minipage}[b]{0.48\textwidth}
%         \centering
%         \begin{tikzpicture}[thick, scale=0.9, baseline=(current bounding box.south)]
%             % unify bbox across BOTH figures
%             \path[use as bounding box] (-0.5,0.5) rectangle (7.5,5.5);

%             \foreach \x in {0,...,7} {
%                 \foreach \y in {2,...,4} {
%                     \node at (\x,\y) [circle, draw=lightgray, fill= lightgray, minimum width=2pt] {};   
%                 }
%             }

%             \draw[line width=0.3mm, black] (0,2) -- (5,3) -- (7,4) -- (2,3) -- cycle;
%             \draw[line width=0.4mm, lightgreen](0,2) -- (7,4);
%             \draw[line width=0.3mm, orange](2,3) -- (5,3);
            

%             \foreach \x in {2,...,5} {
%                 \node at (\x,3) [circle, fill=orange, minimum width=4pt] {};
%             }

%             % \node at (0,2) [circle, fill=black, minimum width=4pt] {};
%             % \node at (7,4) [circle, fill=black, minimum width=4pt] {};

%             \draw(1.1,3) node[black]{$(0,0)$};
%             \draw(5.9,3) node[black]{$(\ell,0)$};
%             \draw(0.5,1.6) node[black]{$(-n,-1)$};
%             \draw(6.5,4.4) node[black]{$(\ell+n,1)$};
%         \end{tikzpicture}
%         \caption{A lattice polygon with fixed \mbox{lattice} diameter (orange) and arbitrarily long \mbox{Euclidean} diameter (green).}
%         \label{fig:unbounded_ed}
%     \end{minipage}
%     \hfill
%     % Right figure
%     \begin{minipage}[b]{0.48\textwidth}
%         \centering
%         \begin{tikzpicture}[thick, scale=0.9, baseline=(current bounding box.south)]
%             % unify bbox across BOTH figures
%             \path[use as bounding box] (-0.5,0.5) rectangle (7.5,5.5);

%             \foreach \x in {0,...,7} {
%                 \foreach \y in {1,...,5} {
%                     \node at (\x,\y) [circle, draw=lightgray, fill= lightgray, minimum width=2pt] {};   
%                 }
%             }
%             \draw[line width=0.3mm, black] (5,1) -- (7,4) -- (2,5) -- (0,2) -- (3,1) -- cycle;

%             % segment
%             \draw[line width=0.3mm, orange](1,3) -- (6,3);

%             \foreach \x in {1,...,6} {
%                 \node at (\x,3) [circle, fill=orange, minimum width=4pt] {};
%             }
%         \end{tikzpicture}
%         \caption{A lattice diameter segment \mbox{in the interior} of a lattice polygon.}
%         \label{fig:LL_interior}
%     \end{minipage}
% \end{figure}
% % \begin{figure}[ht!]
% %     \centering
% %         \begin{tikzpicture}[thick, scale=1]\vspace{0.5cm}
% %         \draw[line width=0.3mm, white](-3,-1.5) -- (6,-1.5);
    
% %         \foreach \x in {-2,...,5} {
%             \foreach \y in {-1,...,1} {
%                 \node at (\x,\y) [circle, draw=lightgray, fill= lightgray, minimum width=1pt] {};   
%             }
%         }
    
%          \draw[line width=0.3mm, black] (-2,-1) -- (3,0) -- (5,1) -- (0,0)-- cycle;
%          \draw[line width=0.3mm, orange](0,0) -- (3,0);
%          \draw[line width=0.3mm, lightgreen](-1.9,-1) -- (4.9,1);

%          \foreach \x in {0,...,3} {
%                 \node at (\x,0) [circle, fill=orange, minimum width=4pt] {};
%             }
%         \node at (-2,-1) [circle, fill=black, minimum width=4pt] {};
%         \node at (5,1) [circle, fill=black, minimum width=4pt] {};

%         \draw(-0.8,0.3)node[black]{$(0,0)$};
%         \draw(4.1,-0.3)node[black]{$(\ell-1,0)$};
%         \draw(-1.6,-1.5)node[black]{$(-n,-1)$};
%         \draw(4.8,1.3)node[black]{$(\ell-1+n,1)$};
        

%     \end{tikzpicture} \vspace{-0.5cm}
    
%     \caption{unbounded Euclidean diameter for fixed lattice diameter}
%     \label{fig:unbounded_ed}
% \end{figure}
% \begin{figure}[ht!]
%     \centering
%     \includegraphics[width=0.8\linewidth]{Figures/1.0_unbounded-ed.png}
%     \caption{unbounded Euclidean diameter for fixed lattice diameter}
%     \label{fig:unbounded_ed}
% \end{figure}

Lattice lines and lattice diameter lines have several interesting connections. For example, in discrete tomography Gardner and Gritzmann have studied the reconstruction of lattice point sets from ``$X$-rays'', i.e., the collection of line sums over parallel lattice lines~\cite{gardnergritzmann97}. Algorithmic connections to combinatorial optimization also exist through the notion of \emph{lattice width}, which is used in algorithms for integer programming. The lattice width has been shown to be bounded by the lattice diameter \cite{Barany_Furedi}. Motivated by questions in the geometry of numbers, the volumes of sections and projections of convex bodies have been studied for many years; see the surveys \cite{giannopoulos2023inequalities, nayar2023extremal}. Analogous questions for the number of lattice points have only been studied more recently \cite{gard05,freyerhenk22,freyerhenk24}. Algorithmically, it is a challenge to determine a section of a convex body with maximal volume, but algorithms exist for convex polytopes~\cite{BDLM}. However, computing a section of a lattice polytope with the most lattice points is still an open problem. In the plane, this problem reduces to finding a lattice diameter segment, which we solve here. 

%Classical inequalities – such as those of Loomis–Whitney~\cite{} and Bourgain’s slicing problem~\cite{} — seek optimal bounds relating the volume of a convex body to the volume of its sections or, dually, 

Furthermore, classifications of convex lattice polytopes have been studied with respect to different invariants. This is of particular interest in algebraic geometry, where reflexive polytopes correspond to specific toric Fano varieties that have ties to mirror symmetry in physics~\cite{batyrev1999classification, kasprzyk2006toric}. Arnold, B\'ar\'any, Pach, Vershik, and others~\cite{arno80, bara92, bara92a} have classified lattice polytopes by their volume while Averkov, Lagarias, Nill, Weismantel, Ziegler, and others \cite{lagariasZiegler1991bounds, averkov2011maximal, nillziegler2011projecting} have continued the classification, based on the number of interior lattice points. We stress that classifying lattice polytopes by their lattice diameter is also possible. A simple result by S. Rabinowitz \cite{rabi89} shows that if $\ldiam(P)<m$, then $|P \cap \Z^d| \leq m^d$.
%, for a lattice polytope $P$.
% \begin{proposition}[\cite{rabi89}]\label{lemma:Rabinowitz}
% Let $P$ be a $d$-dimensional convex lattice polytope and let $m \in \N$ satisfy $|P \cap \Z^d|\geq m^d+1$. Then $P$ has at least $m+1$ collinear lattice points. In other words, if  $\ldiam(P) < m$, then $|P \cap \Z^d| < m^d+1$.
% \end{proposition}
Also in more recent work, Dillon and Arun~\cite{HollowP} asymptotically bound the number of vertices that a lattice polytope with fixed lattice \mbox{diameter can have.}


%For $n=2$ the proof is an elementary pigeon-hole argument. If $|P \cap \Z^d| \geq 2^d$, then there would exist $x,y \in P \cap \Z^d$ whose coordinates all have the same parity, thus the midpoint $\frac{1}{2}x+\frac{1}{2}y \in P \cap \Z^d$ yielding three collinear points in $P$.

% \begin{proof}Consider the coordinates of the integer points modulo $m$.
% Since there are only $m^d$ distinct $d$-tuples of integers modulo
% $m$, some two points ${\bf x}$ and ${\bf y}$ of $P\cap \mathbb{Z}^d$ must be
% congruent (mod $m$). In other words, for all $i=1, 2, \ldots , d$, we have
% $$x_i-y_i\equiv 0\quad ({\rm mod}\ m).$$
% By convexity, all the $m+1$ collinear lattice points
% $${\bf x}+\mbox{${j\over m}$}({\bf y}-{\bf x}),\quad j=0, 1, 2, \ldots
% , m$$ belong to $P$. The lemma is proved.
% \end{proof}

% Furthermore, in the plane Pick's Theorem~\cite{pick99} turns Rabinowitz's bound into an upper bound on the area of $P$ in terms of its lattice diameter. 

Another theme of our paper is Ehrhart theory. In its simplest version, for a $d$-dimensional lattice polytope $P$, and $k \in \N$, the Ehrhart counting function, $\Ehr_P(k) = |kP \cap \Z^d|$, agrees with a polynomial in $k$ of degree $d$. This foundational result extends to rational, multivariate, and weighted versions, using generating functions and featuring reciprocity theorems \cite{beckrobins}. Here we count not individual lattice points, but subsets of them.
It can be shown that there is an Ehrhart theorem for counting lattice lines with fixed direction in a rational polytope, that pass through a lattice point of the rational polytope \cite{aebThesis}.
% It was observed by the second author and collaborators that one has a similar Ehrhart theorem for counting lattice lines in a fixed direction $\bu \in \Z^d$ that contain a lattice point of a rational polytope \cite{unifiedEtheory}. 
In this paper we study the counting function of lattice diameter lines (or segments) of a lattice polygon.

A final motivation for our work is the famous Borsuk's partition problem. In 1933, Borsuk posed the following question in \cite{borsuk1933drei}:
\begin{quote}
    \emph{Can every bounded set $S \subset \R^d$ be partitioned into $d+1$ sets with \mbox{smaller diameter?}}
\end{quote}
Here \textit{diameter} refers to the Euclidean diameter, i.e., $\diam(S) := \max_{\bx,\by \in S} \|\bx-\by\|_2$ where $\| \cdot \|_2$ is the Euclidean norm.
The answer to Borsuk's question is positive for $d=2$, see \cite{borsuk1933drei}, and for $d=3$, see \cite{perkal1947subdivision}, \cite{eggleston1955covering} and \cite{grunbaum1957simple}. In higher dimensions it was believed to be true until Kahn and Kalai \cite{kahn1993counterexample} constructed a counterexample for $d=1325$ and for each $d>2014$. Today, there are known counterexamples in dimensions $64$ and higher, see \cite{jenrich201464}, but the problem is still open for $4\leq d\leq 63$.
More generally, the problem of partitioning a set into sets with smaller diameter is studied for specific classes of bounded sets. In this paper, we study this problem for sets of lattice points with respect to the lattice diameter instead of the Euclidean diameter. We refer to this challenge as the \textit{discrete Borsuk partition problem}.

\subsection*{Our Contributions}

In \Cref{sec:LD-poly} we show that lattice diameter lines of lattice polygons can be computed efficiently:

\begin{manualthm}\emph{\textbf{Theorem~2.4.}}
    Let $P \subset \R^2 $ be a lattice polygon. There exists a polynomial-time algorithm that computes a lattice diameter line of $P$ in every direction in which such a line exists. In particular, the lattice diameter of $P$ can be computed in polynomial time.
\end{manualthm} 

In contrast, we show the following hardness result in \Cref{sec:NP-hard}:

\begin{manualthm}\textbf{\emph{Theorem~2.7.}}
    Let $d \geq 3$ and let $K \subset \R^d$ be a $d$-dimensional bounded semi-algebraic set. \mbox{Computing} a lattice diameter of $K$ is an $\NP$-hard problem. 
    %Moreover, even computing a $\bu$-lattice diameter of $K$ in a fixed direction $\bu \in \Z^d \setminus \{0\}$ is $NP$-hard.
\end{manualthm}

In \Cref{sec:QP} we generalize Ehrhart theory to counting the set of lattice diameter lines. While counting the number of lattice lines in a fixed direction that intersect $P$ in a lattice point is not hard \cite{aebThesis}, counting those that are lattice diameter lines is much more difficult. Surprisingly, we have an Ehrhart-type result.

\begin{definition}\label{def:LD_P}
    For a lattice polytope $P$, let $\LD_P(k)$ denote the number of lattice diameter lines of the dilated polytope $kP$.
\end{definition}

\begin{manualthm}
    \emph{\textbf{Theorem~3.6.}} Let $P \subset \R^2$ be a lattice polygon. There exists a positive integer $q$ such that \mbox{$\LD_P: \N_{\geq q} \to \N$} agrees with a quasi-polynomial function of degree one whose period divides $q$.
\end{manualthm}

 In \Cref{sec:directions-Borsuk} we begin by generalizing two results by Alarcon~\cite{Alarcon_PhD} to higher dimensions. The first result bounds the total number of lattice diameter directions:

\begin{manualthm}\textbf{\emph{Theorem~4.3.}}
    Let $P\subset \R^d$ be a lattice $d$-polytope with $\ldiam(P)=1$. Then $P$ has at most $\binom{2^d}{2}$ lattice diameter directions and this bound is best possible.
\end{manualthm}

The second generalization is about a local property of lattice diameter segments. In fact, we prove it more generally for bounded sets. This result is a key ingredient to prove the discrete version of Borsuk's partition problem.

\begin{manualthm}\textbf{\emph{Lemma~4.6.}}
    Let $S\subset \Z^d$ be a bounded set and let $\bx \in S$. Then, the number of lattice diameter segments of $S$ that have $\bx$ as an endpoint is bounded by $2^d-1$, and this bound is best possible.
\end{manualthm}

To study the discrete version of Borsuk's partition problem, we define the following notion:

\begin{definition}
    Let $S \subset \Z^d$ be a bounded set. Define the \textit{lattice Borsuk number} of $S$, $\beta_\Z(S)$, as the minimal size of a partition of $S$ where each part has a smaller lattice diameter than $\ldiam(S)$.
\end{definition}

Finally, the following theorem answers the discrete version of Borsuk's partition problem:

\begin{manualthm}\textbf{\emph{Theorem~4.8.}}
    Let $S \subset \Z^d$ be a bounded set. Then, $\beta_\Z(S) \leq 2^d$ and this bound is best possible.
\end{manualthm}

In particular, \Cref{Thm:Borsuk} says that $\beta_\Z(P \cap \Z^d) \leq 2^d$ for $d$-dimensional polytopes $P$. Note that we are not partitioning $P$ into smaller polytopes, i.e., the parts of $P \cap \Z^d$ are not necessarily the lattice points of a polytope. The bound $2^d$ is also best possible for polytopes; it is attained for the lattice points of a $d$-dimensional cube.

\subsection*{Notation}
We use basic notions and notation as can be found in the textbooks~\cite{grub07,zieglerbook}.
% A convex body means a compact convex set, and a lattice polytope in $\R^d$ is the convex hull of finitely many integer points in $\Z^d$.
%Vectors $\bv \in \R^d$ will be written bold.
A $d$-polytope $P$ is the convex hull of finitely many points in $\R^d$ whose affine hull is $d$-dimensional. It is called a \emph{lattice} (or \emph{rational}) $d$-polytope if its vertices belong to the integer lattice $\Z^d$ (or to $\Q^d)$. We say two lattice $d$-polytopes $P$ and $Q$ are unimodularly equivalent if there exists a unimodular matrix $A \in \GL(d,\Z)$ and $\bt \in \Z^d$ such that $P=AQ+\bt$, i.e., these transformations correspond to lattice preserving maps. We write $\spn(S), \, \aff(S)$, and $\conv(S)$ for the linear, affine, and convex hull of a set $S \subset \R^d$, respectively.
The segment between two points $\bx,\by \in \R^d$ is denoted $[\bx,\by]:=\conv(\{\bx,\by\})$ and when $\bx,\by \in \Z^d$, we call this a \textit{lattice segment}. A \emph{line} will always mean an \emph{affine line} and a \emph{lattice line} is a line that passes through at least two lattice points. We say that a line has direction $\bu$ if it can be written as $\bx+\spn(\bu)$ for some $\bx \in \R^d$. We say $\bu$-lattice line (or segment) if the lattice line (or segment) has direction $\bu$. Furthermore, from now on a \textit{diameter direction} means a lattice diameter direction. The inner product we use is the standard inner product on $\R^d$, denoted $\langle \cdot, \cdot \rangle$. A hyperplane and its closed half-spaces are denoted $H(\ba,\alpha):= \{ \bx \in \R^d \st \langle \ba, \bx \rangle = \alpha \}$, $H^+(\ba,\alpha)$ and $H^{-}(\ba,\alpha)$, respectively. For a Lebesgue measurable set $S\subset \R^d$, we write $\vol(S)$ which means the intrinsic volume of $S$ in $\aff(S)$. Furthermore, a vector $\bu=(u_1, \ldots, u_d) \in \Z^d$ is called primitive if $\gcd(u_1, \ldots, u_d) =1$. For a positive integer $n$ we write $[n]:=\{1, \ldots, n\}$ and $\binom{[n]}{k}$ denotes the $k$-element subsets of $[n]$. For integers $a,q \in \Z$ we write $a \bmod q$ to mean the representative of $a$ in $\{0, \ldots, q-1\}$.
%Lastly, the natural numbers $\N$ contain $0$.
%, and when we write $d \in \N$ for dimension we assume $d\geq 2$.
In \Cref{sec:directions-Borsuk}, we use graph theory and standard notation as in \cite{bondy2008graph}. A graph is a simple undirected graph $G=(V,E)$, where $V$ is a finite set and $E \subseteq \binom{V}{2}$. In this context, $V$ or $V(G)$ is the set of \emph{vertices} and $E$ or $E(G)$ is the set of \emph{edges} of the graph $G$. Two vertices $u,v\in V$ are called \emph{adjacent} if $\{u,v\} \in E$.

%OLD introduction notes

% \begin{definition}\label{def:ld-convex}
% Let $C \subset \R^d$ be a convex body. An affine lattice line $L$ is called a \textit{lattice diameter line} of $C$ if $|L \cap C \cap \Z^d| \geq | L' \cap C \cap \Z^d|$ for all affine lines $L'$. In this case $L \cap C$ is called a \emph{lattice diameter segment}, and $\ldiam(C):=|L \cap C \cap \Z^d|-1$ is called the \emph{lattice diameter} of $C$. Further, for a non-zero vector $\bu \in \Z^d$, a \textit{$\bu$-lattice line} $L$ is an affine lattice line $L$ with direction $\bu$. It is called a \textit{$\bu$-lattice diameter line} of $C$, if $L \cap C \cap \Z^d| \geq |L' \cap C \cap \Z^d|$ for all $\bu$-lattice lines $L'$ and in this case, $|L \cap C|$ is called a $\bu$-\emph{lattice diameter segment} and $|L \cap C \cap \Z^d|-1$ is called the $\bu$\textit{-lattice diameter} of $C$. If a $\bu$-lattice diameter line is a lattice diameter line of $C$, then $\bu$ is called a \textit{diameter direction} of $C$.
% \end{definition}

% counted the number of lattice polytopes with a fixed number of interior points and provide explicit upper bounds of the lattice diameter in terms of the dimension and the number of interior points. What happens when there are no lattice points in the interior? Lattice polytopes without interior points are of interest but their classification is incomplete. Nill and Ziegler~\cite{nillziegler2011projecting} showed that up to unimodular equivalence there are only finitely many $d$-dimensional lattice polytopes without interior lattice points that do not admit a lattice projection onto a ($d-1$)-dimensional lattice polytope without interior lattice points. This connects to work of  Averkov, Wagner and Weismantel~\cite{AverkovWW10}, on the finiteness of the number of maximal lattice polytopes without any interior lattice points. 

% Furthermore, convex lattice polytopes with bounded volume or few interior points have been classified in various settings~\cite{Arn80, BP92, BV92, LZ91, NZ11} [CHECK THESE CITATIONS!]. 
\section{Computing lattice diameters and lattice diameter lines}\label{sec:LD-computation}

In this section we will explain and prove \Cref{thm:LD-poly} and \Cref{thm:NP-hard}.

\subsection{A polynomial-time algorithm to compute lattice diameter lines in the plane}\label{sec:LD-poly}

Let $P$ be a lattice polygon. We present an algorithm that runs in polynomial time in the input size of $P$ and computes a lattice diameter line of $P$ for every direction in which one exists. 
% The input size of the algorithm is $n\cdot \log_2(c)$ where $n$ is the number of vertices in $P$ and $c$ is the largest absolute value of a coordinate of a vertex. Since the edges of $P$ can be obtained from the vertices in polynomial time, we may assume that the edges of $P$ also belong to the input. We reduce the problem to finding \textit{local lattice diameter lines}. A \textit{local lattice diameter line} at vertex $\bv$ of a lattice polygon $P$, is an affine line $L$ intersecting $\bv$ such that $|L \cap P \cap \Z^2|$ is maximal among all affine lines intersecting $\bv$. See \Cref{fig:local-ld}.
% %\input{Figures/2.0_local-ld}
% \begin{figure}[ht!]
%     \centering
%     \includegraphics[width=0.75\linewidth]{Figures/2.0_local-ld.png}
%     \caption{}
%     \label{fig:local-ld}
% \end{figure}
The algorithm relies on the following three claims:
\begin{itemize}
    \item For every diameter direction $\bu$, there exists a lattice diameter line in direction $\bu$ passing through a vertex $\bv$ and an edge $e$ such that some translate of $\aff(e)$ is tangent to $P$ at $\bv$. In this case, we say that $\bv$ is an \textit{opposite vertex} to the edge $e$. %This is \Cref{lemma:vertex_edge_pair}.
    \item For every pair $(e,\bv)$, where $e$ is an edge and $\bv$ is an opposite vertex to $e$, one can compute in polynomial time a \textit{local lattice diameter line} at $\bv$ of the triangle $\conv \{ e, \bv\}$, i.e., a line through $\bv$ that intersects $\conv(\{e,\bv\})$ in the most lattice points.
    %This relies \Cref{lemma:min-gens}
    \item At most three lattice diameter lines pass through any fixed vertex. This was proven by Alarcon \cite[Lemma 2.4]{Alarcon_PhD} (and it also follows from our \Cref{Lemma:MaxDegree}).
\end{itemize}

The algorithm is as follows:

\begin{algorithm}[ht!]
\caption{\textsc{Computing Lattice Diameter Lines}}
\textbf{Input:} a lattice polygon $P$ given by its vertices \\
\textbf{Output:} at least one lattice diameter line for each direction in which a lattice diameter line exists
\begin{algorithmic}[1]
\State Initialize two empty sets: $\Tcal$ for triangles, and $\Lcal$ for candidate lines
\ForAll{edges $e$ of $P$}
    \State compute all opposite vertices $\bv$ to $e$ (each edge has at most two opposite vertices)
    \State form triangles $\conv(\{e,\bv\})$ and add them to $\Tcal$
\EndFor
\ForAll{triangles in $\Tcal$}
    \State compute up to three local lattice diameter lines of $T=\conv(\{e,\bv\})$ at $\bv$
    \State add these local lattice diameter lines to $\Lcal$
\EndFor
\State \Return all lines in $L \in \Lcal$ that maximize $|L \cap P \cap \Z^2|$
\end{algorithmic}
\end{algorithm}

We proceed to prove some lemmas. Recall that for a primitive vector $\bu \in \Z^d\setminus \{0\}$, and a $\bu$-lattice line $L$, the lattice normalized length of the segment $L \cap P$ is $\nvol(L \cap P) := \frac{\vol(L \cap P)}{\|\bu\|}$.

\begin{lemma}\label{lem:nvol}
    Let $P$ be a lattice polytope and $L$ be a line that intersects $P$. Then
    \begin{align*}
    \left \lfloor \nvol(L \cap P) \right \rfloor \leq |L \cap P \cap \Z^d| \leq  \left \lfloor \nvol(L \cap P) \right \rfloor + 1
    \end{align*}
and the upper bound is attained if $L$ contains a lattice point on the boundary of $P$.
\end{lemma}

%Anouks Proof to wordy/long
% \begin{proof}
%     Let $L=(x+\spn\{u\})$ where $u \in \Z^d$ is primitive. Write $L \cap P=[a,b]$ and let $a',b' \in [a,b] \cap \Z^2$ be the closest lattice points to $a$ and $b$, respectively. Then
%     \begin{align}\label{eq:line_ineq}
%         \left( \big | [a,b]\cap \Z^2 \big |-1 \right) \cdot \|\bu\|= \left( \big | [a',b']\cap \Z^2 \big |-1 \right) \cdot \|\bu\| = \vol([a',b'])
%     \end{align}
%     as depicted in \Cref{fig:bound_nlp}.
    
%     %\input{Figures/2.12_bound_nlp}
%     \begin{figure}[ht!]
%     \centering
%     \includegraphics[width=0.66\linewidth]{Figures/2.12_bound_nlp.png}
%     \caption{number of lattice points on a line in comparison to its Euclidean length}
%     \label{fig:bound_nlp}
%     \end{figure}

    
%     Plugging $ \vol( [a',b'] )\leq \vol( [a,b] )$ into \eqref{eq:line_ineq}
%     we get $\left( \big | [a,b]\cap \Z^2 \big |-1 \right) \cdot \|\bu\| \leq \vol([a,b])$, and the right-hand-side follows from rearranging this inequality, and using the fact that $| [a,b]\cap \Z^2|$ is an integer.
%     Also $\vol([a,b])-\vol([a',b']) < 2 \|\bu\|$
%     which together with \eqref{eq:line_ineq} yields
%     \begin{align}\label{eq:nvol-LHS}
%         \nvol([a,b]) -1 < |[a,b] \cap \Z^2|.
%     \end{align} If $\nvol([a,b]) \in \Z$ then the strict inequality in \eqref{eq:nvol-LHS} implies $\nvol([a,b]) \leq |[a,b] \cap \Z^2|$. If $\nvol([a,b]) \notin \Z$, then $ \left \lfloor \nvol([a,b])\right \rfloor = 
%     \left \lceil \nvol([a,b]) \right \rceil -1 \leq |[a,b] \cap \Z^2|$. This proves the left-hand-side. For the equality: Suppose $a=a'$ is a lattice point, then $\vol([a,b])=\vol([a',b'])+\vol([b',b])$, and $\vol([b',b])< \|\bu\|$ which implies that 
%     \begin{align*}
%         \bigg \lfloor \frac{\vol([a,b])}{\|\bu\|} \bigg \rfloor = \bigg \lfloor  \underbrace{\frac{\vol([a',b'])}{\|\bu\|}}_{\in \N} + \underbrace{\frac{\vol([b',b])}{\|\bu\|}}_{\in [0,1)} \bigg \rfloor = \frac{\vol([a',b'])}{\|\bu\|} = \big |[a',b'] \cap \Z^2 \big | -1 = \big |[a,b] \cap \Z^2 \big | -1.
%     \end{align*}
% \end{proof}

\begin{proof}
Parametrize the line as $L = \{\bx_0 + t\bu : t \in \mathbb{R}\}$, where $\bx_0$ is an arbitrary lattice point on $L$ and $\bu \in \mathbb{Z}^d$ is a primitive vector in the direction of $L$. Then, the intersection $L \cap P$ is a segment with endpoints $\ba = \bx_0 + t_1\bu$ and $\bb = \bx_0 + t_2\bu$, for some real numbers $t_1 \leq t_2$.
The lattice points in $L \cap P$ are precisely those points of the form $\bx_0 + k\bu$ with $k \in \mathbb{Z} \cap [t_1, t_2]$. Moreover,
\[
\nvol(L \cap P) = \frac{\|\bb - \ba\|}{\|\bu\|} = \frac{\|(t_2 - t_1) \bu\|}{\|\bu\|} = |t_2 - t_1|=t_2-t_1.
\]
% Using $\lceil t_1 \rceil - 1 < t_1 \leq \lceil t_1 \rceil$ and $\lfloor t_2 \rfloor \leq t_2 < \lfloor t_2 \rfloor + 1$ we obtain

If $t_1$ is not an integer, then $|L \cap P \cap \mathbb{Z}^d| 
= \lfloor t_2 \rfloor - \lceil t_1 \rceil + 1 = \lfloor t_2 \rfloor - \lfloor t_1 \rfloor$
and the claim follows from the observation
$\lfloor t_2-t_1 \rfloor \leq \lfloor t_2 \rfloor - \lfloor t_1 \rfloor \leq \lfloor t_2-t_1 \rfloor +1$. If at least one of $t_1$ or $t_2$ is an integer, then
$|L \cap P \cap \Z^d| = \lfloor t_2 \rfloor - \lceil t_1 \rceil + 1 = \lfloor t_2 - t_1 \rfloor + 1=\lfloor \nvol(L \cap P) \rfloor +1$.
\end{proof}

%TRIANGLE INEQUALITY FOR FLOORING 
% \begin{proposition}
%     Let $a,b \in \R$. Then $\lfloor a \rfloor + \lfloor b \rfloor \leq \lfloor a+b \rfloor$. Equivalently $\lfloor b \rfloor - \lfloor a \rfloor \geq \lfloor b-a \rfloor$.
% \end{proposition}

% \begin{proof}
%     The first inequality should be clear. For a proof, it suffices to consider $a,b \in (0,1)$ (after subtracting the integer part of both numbers) and then the equality reduces to $0 \leq 0$ or $0 \leq 1$. The equivalence to the second inequality comes from letting $c=a+b$ and rearranging.
% \end{proof}

%MORE EXPLANATIONS
% \begin{proposition}
% Let $t_1 < t_2$ be two real numbers. Then
% \[
% \lfloor t_2 \rfloor - \lceil t_1 \rceil + 1 \geq \lfloor t_2 - t_1 \rfloor.
% \]
% \end{proposition}

% \begin{proof}
% We have $\lceil t_1 \rceil - 1 < t_1 \leq \lceil t_1 \rceil$ and $\lfloor t_2 \rfloor \leq t_2 < \lfloor t_2 \rfloor + 1$.

% The number of integers in the interval $[t_1, t_2]$ is
% \[
% N := \lfloor t_2 \rfloor - \lceil t_1 \rceil + 1.
% \]

% Observe that
% \[
% t_2 - t_1 < (\lfloor t_2 \rfloor + 1) - (\lceil t_1 \rceil - 1) = \lfloor t_2 \rfloor - \lceil t_1 \rceil + 2,
% \]
% and also
% \[
% t_2 - t_1 > \lfloor t_2 \rfloor - \lceil t_1 \rceil.
% \]

% Therefore,
% \[
% \lfloor t_2 \rfloor - \lceil t_1 \rceil < t_2 - t_1 < \lfloor t_2 \rfloor - \lceil t_1 \rceil + 2,
% \]
% so the floor of $t_2 - t_1$ is at most $\lfloor t_2 \rfloor - \lceil t_1 \rceil + 1$, that is,
% \[
% \lfloor t_2 - t_1 \rfloor \leq \lfloor t_2 \rfloor - \lceil t_1 \rceil + 1 = N.
% \]

% This concludes the proof.
% \end{proof}

As mentioned above, we say that a vertex $\bv$ of a lattice polygon $P$ is an \emph{opposite vertex} to an edge $e$ of $P$ if the furthest translate of $\aff(e)$ that intersects $P$ passes through $\bv$. More formally, if $\ba$ is an outward normal vector of the edge $e$, then $\langle \ba, \bv \rangle = \min_{\bx \in P} \langle \ba,\bx \rangle$. See \Cref{fig:edge-far-vertex}.
%EXAMPLE: Lemma 2.3 <v,a> = min <x,a> picture

\begin{figure}[H]
    \centering
        \begin{tikzpicture}[thick, scale=0.8]\vspace{0.5cm}
        \draw[line width=0.3mm, white](0,0.5) -- (7,0.5) ;
    
        \foreach \x in {-2,...,8} {
            \foreach \y in {0,...,5} {
                \node at (\x,\y) [circle, draw=lightgray, fill= lightgray, minimum width=1.5pt] {};   
            }
        }
    
         \draw[line width=0.3mm, darklav, opacity=0.8](3.5,0) --+ (2.5,5);
         \draw[line width=0.3mm, darklav, opacity=0.8](-1,0) --+ (2.5,5);
         \draw[line width=0.3mm, black] (1,1) -- (0,2)-- (3,4) -- (4,4)  -- (5,3)  -- (4,1) -- cycle;
         \draw[line width=0.6mm, black] (5,3)  -- (4,1);
         \draw[line width=0.3mm, orange](0,2) -- (4.5,2);
         %\draw[line width=0.3mm, darklav,-{Latex[round]}](4.75,2.5) --+ (0.8,-0.4);

         
         \draw[line width=0.3mm, darklav,-{Latex[round]}](6,1) --+ (0.8,-0.4);
         \node at (6,1) [circle, fill=black, minimum width=4pt] {};
         \draw(7,0.5)node[darklav]{$\ba$};
         \draw(5.8,0.5)node[black]{$(0,0)$};
         
         \foreach \p in {(0,2), (1,2),(2,2),(3,2), (4,2)} {
                \node at \p [circle, fill=orange, minimum width=4pt] {};
            }
       
       \draw(-0.5,2)node[orange]{$\bv$};
       \draw(2.35,2.37)node[orange]{$L \cap P$};
       \draw(5.5,2.4)node[black]{$\footnotesize \text{edge } e$};
       \draw(-2.2,2)node[orange]{$\footnotesize \text{opposite vertex} $};
       \draw(-0.7,3.5)node[darklav]{$H(\ba,\langle \ba, \bv \rangle)$};
       \draw(6,3.5)node[darklav]{$\footnotesize\aff(e)$};
 \end{tikzpicture} 
    \caption{A lattice diameter line passing through an edge and its opposite vertex.}
    \label{fig:edge-far-vertex}
\end{figure}

If $\bu \in \Z^2 \setminus \{0\}$, then we call a $\bu$-lattice line $L$ a \emph{$\bu$-lattice diameter line} of $P$, if among all lines in direction $\bu$, it contains the most lattice points in $P$. The following lemma shows that a $\bu$-lattice diameter line of $P$ can always be chosen to pass through an edge of $P$ and an opposite vertex. This follows from \cite[Lemma~1]{Barany_Furedi}, where it is stated without a proof. We include one here for completeness.

\begin{lemma}\cite[Lemma 1]{Barany_Furedi}\label{lemma:vertex_edge_pair}
Let $P$ be a lattice polygon and $\bu \in \Z^2 \setminus \{0\}$. There exists a $\bu$-lattice diameter line $L$ that passes through a vertex $\bv$ of $P$ and through an edge $e$ of $P$ with outward normal vector $\ba$ such that $\langle \ba, \bv \rangle = \min_{\bx \in P} \langle \ba,\bx \rangle$.
\end{lemma}

% A lattice diameter line can also pass through two vertices. In this case \Cref{lemma:vertex_edge_pair} still holds. It says that for one of the two vertices (call it $\bv$) and one of the edges that the other vertex is contained in (call it $e$), $\bv$ is an opposite vertex to $e$.

\begin{proof}
    \Cref{lem:nvol} implies that if a $\bu$-lattice line $L$ passes through a vertex of $P$, and satisfies $\vol(L\cap P) \geq \vol(L' \cap P)$ for all $\bu$-lattice lines $L'$, then $L$ is a $\bu$-lattice diameter line of $P$.

    First we prove that such a line exists. Let $L$ be \textit{any} line in direction $\bu$ satisfying $\vol(L\cap P) \geq \vol(L' \cap P)$ for all lines $L'$ in direction $\bu$ (by compactness of $P$ such a line exists). If $L$ passes through parallel edges, then a translate of $L$ along $\bu^{\perp}$ will pass through a vertex of $P$ and its intersection with $P$ will have the same length. Else, translating $L$ in one direction along $\bu^{\perp}$, until it passes through a vertex, will strictly increase the length of its intersection with $P$, a contradiction. This proves the existence of such a line, and necessarily, it will be a $\bu$-lattice line.
    
    So let $L$ be a $\bu$-lattice line that passes through some vertex of $P$ and satisfies
    $\vol(L\cap P) \geq \vol(L' \cap P)$ for all $\bu$-lattice lines $L'$.
    First, suppose that $L$ passes through the relative interior of an edge $e$ and through a vertex $\bv$ of $P$. Towards a contradiction, suppose $\langle \ba, \bv \rangle \neq \min_{\bx \in P} \langle \ba,\bx \rangle$. Let $\tilde{\bv}$ be a vertex for which $\min_{\bx \in P} \langle \ba,\bx \rangle$ is attained. It is an elementary trigonometry calculation to show that there exists a point $\bp \in [\bv,\tilde{\bv}] \subseteq P$ such that a translate of $L$ passing through $\bp$ also passes through $e$, and its intersection with $P$ is longer.
    For the second case, suppose $L$ passes through two vertices $\bv$ and $\bw$. Using convexity of $P$ and the fact that $L \cap P$ is the longest segment in $P$ in direction $\bu$ one can find parallel supporting hyperplanes (lines) of $P$ at $\bv$ and at $\bw$. Rotating them, until one of the lines becomes tangent to an edge of $P$ proves this case.
\end{proof}

%\input{Figures/2.13_p-longer}

%EXPLANATION:
% Fixing $\alpha$, the distance (in direction $a$) of a parallel hyperplane $H(a,\beta)$ to $H(a,\alpha)$ where $\alpha>\beta$, is given by $\delta_a:=\frac{\alpha-\beta}{\|a\|}$, thus the distance increases if $\beta$ decreases. Note that $\langle u,\ba \rangle \neq 0$ since $L$ has direction $\bu$ and passes through the interior of $P$, and an edge with normal vector $a$. Thus the distance between parallel hyperplanes $H(a,\alpha)$ and $H(a,\beta)$ in direction $\bu$ is given by $\delta_u:=\delta_a \cdot \frac{\langle u,\ba \rangle}{\|\bu\|\|a\|}=\frac{\alpha-\beta}{\|a\|} \cdot \frac{\langle u,\ba \rangle}{\|\bu\|\|a\|}$ (this is jus ttrigonometry)
% %(This is trigonometry: Consider the right angled triangle with hypotenuse side $(x+\spn\{u\}) \cap P$, adjacent side $(x+\spn\{a\}) \cap P$, and the opposite side contained in $H(a,\alpha)$. Let $\theta$ be the angle between the adjacent and hypotenuse, see \Cref{fig:dist-hp}. It is $\delta_a/\delta_u = \cos \theta$. Also $\cos \theta = \frac{\langle u,a\rangle}{\|\bu\|\|a\|}$, so $\delta_a = \delta_u \cdot \frac{\langle u,\ba \rangle}{\| u \| \|a\|}$.), 
% so the distance in direction $\bu$ increases if and only if $\beta$ decreases.
% %\input{Figures/2.13_dist-hp} 
% \begin{figure}[ht!]
%     \centering
%     \includegraphics[width=0.61\linewidth]{Figures/2.13_dist-hp.png}
%     \caption{distance of two hyperplanes in direction $a$ and direction $\bu$}
%     \label{fig:dist-hp}
% \end{figure}


% Since $L$ passes through the interior of $e$, there exists a small translation of $L$ to $L'$ such that $L'$ passes through a point $p \in [v,\tilde{v}] \subseteq P$ as well as the same edge $e \subseteq H(a,\alpha)$, and thus  $\vol(L'\cap P)$ is exactly the distance (in direction $\bu$) of the hyperplanes $H(a,\alpha)$ and $H(a,\langle p,a\rangle)$. Since $\beta:= \langle p, a \rangle < \langle v, a \rangle$, the Euclidean length of $L' \cap P$ is longer than the Euclidean length $L \cap P$, a contradiction.


%LONG VERSION
% \textit{Case 2:} Suppose $L$ passes through two vertices $v_L$ and $v_R$.  In \Cref{fig:hyp-arr} we assume that the diameter direction is $u=(1,0)$, which one can always achieve after applying a unimodular transformation to $P$.
% % This explains the notation of subscripts ``$L$'' for left and ``$R$'' for right.
% For $i=1,2$ let $v_L + H_{L_i}$ and $v_R + H_{R_i}$ be the support hyperplanes of the two edges incident to $v_L$, and $v_R$, respectively, see \Cref{fig:hyp-arr}.
% %\input{Figures/2.13_hyp-arr}
% \begin{figure}[ht!]
%     \centering
%     \includegraphics[width=0.98\linewidth]{Figures/2.13_hyp-arr.png}
%     \caption{hyperplane arrangement corresponding to the four edges at $v_L$ and $v_R$}
%     \label{fig:hyp-arr}
% \end{figure}

% Define the positive and negative halfspaces of $H_{L_i}$ and $H_{R_i}$ such that $P \subseteq v_L +H^{-}_{L_i}, \ v_R +H^{-}_{R_i} $. We consider the ``$L$''- and ``$R$''-hyperplane arrangements, defined the $H_{L_i}$ and $H_{r_i}$ for $i=1,2$ respectively. Define $L(\sgn_1, \sgn_2):= H^{\sgn_1}_{L_1} \cap H^{\sgn_2}_{L_2}$ for $\sgn_1,\sgn_2 \in \{+,-\}$, and $R(\sgn_1, \sgn_2)$ analogously, as the four cones of the hyperplane arrangement, that partitions $\R^2$ (up to lower-dimensional intersection). So $v_L + L(-,-)$, $v_R + R(-,-)$, are the tangent cones of $P$ at $v_L$ and $v_R$, respectively.

% An edge $e_{L_i}$ is tangent to $P$ at $v_R$ if and only if $H_{L_i}$ is tangent the the cone $R(-,-)$ if and only if $H_{L_i} \subseteq R(+,-) \cup R(-,+)$, and analogously when swapping $L$ with $R$.

% Suppose $H_{L_1},H_{L_2}$ are both not tangent to $R(-,-)$, that is, $H_{L_1},H_{L_2} 
% \subseteq R(+,+) \cup R(-,-)$. Since the diameter direction $\bu$ points inside the tangent cones of $P$ at $v_R$ and $v_L$, it follows that $ L(+,+) \cup L(-,-) \subseteq R(+,+) \cup R(-,-)$, and so both $H_{R_i} \subseteq L(+,-) \cup L(-,+)$ which proves the claim.


%TOÑO IDEA
% \Tonosays{For \Cref{lemma:vertex_edge_pair}, I don't think we need to be too detailed, since it's not our result. Here's a shorter proof that skips the technicalities: \hskip .5cm
% \begin{proof}
% The equality $\langle v, a \rangle = \min_{x \in P} \langle x,\ba \rangle$ means that the translate of $\aff(e)$ passing through $\bv$ is tangent to $P$; that is, it is the furthest translate of $\aff(e)$ in direction $-a$ that still intersects $P$.\\
% Let $L_1,L_2$ be two supporting lines of $P$ orthogonal to $\bu$. These lines intersect $P$ either along an edge or at a single vertex. If one of the lines, say $L_1$, intersects $P$ along an edge $e$, we can take any vertex $\bv$ of $P$ lying on the other line $L_2$. Then the translate of $\aff(e)$ passing through $\bv$ lies in $L_2$ and is therefore tangent to $P$.\\
% If both lines intersect $P$ only at a single vertex, we can rotate them continuously (either clockwise or counterclockwise) around their respective contact points with $P$, keeping $P$ enclosed between them at all times. Eventually, one of the lines—say $L_1$—will intersect $P$ along an edge $e$, while the other, $L_2$, still contains at least one vertex $\bv$ of $P$. This reduces the situation to the previous case, completing the proof.
% \end{proof}}

% \Anouksays{In proof we still need to show that the line $L$ with longest length in direction $\bu$ passes through vertex and edge that are tangent as described. These lines don't have to be orthogonal to the direction $\bu$.}

%OLD PROOF (TOO LONG!)
% Without loss of generality\footnote{Else apply an affine unimodular transformation; these preserve all lattice lines.} we can assume that $u=(1,0)$, one of the vertices is $v=(0,0)$, and the other vertex is $w=(w_1,0)$ with $w_1>0$. One can draw all cases of hyperplane angles and conclude that the claim is true. We abstract these ideas to show that the claim.

% Let $e_L,d_L$ ($L$ for ``left'') be the two edges adjacent to $v=(0,0)$ and $e_R,d_R$ ($R$ for ``right'') be the two edges adjacent to $w$, where $e_L,e_R$ are in the upper half-plane and $d_L,d_R$ are in the lower half-plane. Write $H(e) := \aff(e)$ which is the supporting hyperplane, in this case a line, of an edge $e$. We characterize these lines by their angle from the positive $\bx$-axis, that is, if $H(e)=\spn\{z\}$, for $e \in \{e_R,e_L,d_R,d_L\}$ then we write the angle as $a(e):=\min\{ \angle (0,z), \angle (0,-z)\}$, see \Cref{fig:VV_setup}.
% \begin{figure}[h!]
%     \begin{center}
%     %\includegraphics[width=0.8\linewidth]{Figures/VV-setup.jpeg}
%     \input{Figures/2.13_VV-setup.tikz}
%     \caption{parametrizing hyperplanes of vertices}
%     \label{fig:VV_setup}
%     \end{center}
% \end{figure}

% Now by our assumptions, we get:
% \begin{enumerate}
%     \item $a(e_L) \leq a(e_R)$ and $a(d_R) \leq a(d_L)$, otherwise, translating $\spn\{u\}$ parallel upwards (or downwards) would yield a line in direction $\bu$ of greater Euclidean length, which is a contradiction to $\spn\{u\}$ through $\bv$ and $w$ having maximal Euclidean length, see \Cref{fig:VV-conditions}.
%     \item $a(d_R) < a(e_R)$ and $a(e_L) < a(d_L)$. This follows from convexity of the polygon since $P \subseteq H^+(e_L)$ and $d_L \subseteq H^-(u,0)$, see \Cref{fig:VV-conditions}.
%     \item Combining 1.~and 2., we get $\max \{a(e_L), a(d_R) \}\leq \min \{a(e_R), a(d_L)\}$. This means that, for at least one of the two sides (left or right), the angle of one edge lies between the angles of the two edges on the opposite side. This condition is equivalent to the translated hyperplane being tangent to the other side's vertex. More precisely: without loss of generality, suppose $a(e_L) \leq a(e_R) \leq a(d_L)$. So $H(e_R)$ separates the hyperplanes $H(d_L)$ and $H(e_L)$. More generally, a line through the origin with angle $a$ is tangent to $P$ at $v=(0,0)$ if and only if it lies in $H^-(e_L) \cap H^-(d_L)$. Lying in $H^-(e_L)$ means that the line is left of $e_L$ in the upper half-plane, so $a\geq a(e_L)$. Lying in $H^-(d_L)$ means that the line is left of $d_L$ in the lower half-plane, so $a+180 \leq a(d_L)+180$, or equivalently $a \leq a(d_L)$. Since $a(e_L) \leq a(e_R) \leq a(d_L)$, it follows that $H(e_R)$ is tangent to $P$ at $(0,0)$.
% \end{enumerate}
% \begin{figure}[h!]
%     \centering
%      \input{Figures/2.13_VV-conditions.tikz}
%      \vskip -1.6cm
%     \caption{conditions on edge angles}\label{fig:VV-conditions}
% \end{figure}

Now consider all triangles $T_e=\conv(\{e,\bv_e\})$ where $e$ is an edge of $P$, and $\bv_e$ is an opposite vertex to $e$. \Cref{lemma:vertex_edge_pair} implies that for every diameter direction $\bu$ there exists a lattice diameter line that is a local lattice diameter line of $T_e$ at $\bv_e$ for some edge $e$. \Cref{lemma:min-gens} is the key ingredient to compute local lattice diameter lines of lattice triangles.

\begin{lemma}\label{lemma:min-gens}
    Let $T = \conv \{e,\bv\}$ be a lattice triangle, where $e$ is an edge of $T$ with outward normal vector $\ba$, and $\bv$ is its opposite vertex. If $\bw_1,\bw_2 \in T \cap \Z^2 \setminus \{\bv\}$ satisfy $\langle \ba, \bw_1 \rangle \leq \langle \ba, \bw_2 \rangle$, then for $L_1:= \aff(\{\bv,\bw_1\})$ and $L_2:= \aff(\{\bv,\bw_2\})$
    \[
    |L_1 \cap T \cap \Z^2| \geq |L_2 \cap T \cap \Z^2|.
    \]
\end{lemma}

\begin{proof}
    If $\langle \ba, \bw_1 \rangle = \langle \ba, \bw_2 \rangle$, then $\bw_1$ and $\bw_2$ lie on the same hyperplane $H(\ba, \langle \ba, \bw_1 \rangle)$, and by  proportionality of lines 
\[
\frac{\operatorname{vol}(L_1 \cap T)}{\|\bw_1- \bv\|} = \frac{\operatorname{vol}(L_2 \cap T)}{\|\bw_2- \bv\|}.
\]
Applying \Cref{lem:nvol}, we obtain
\[
|L_1 \cap T \cap \mathbb{Z}^2| = \left\lfloor \frac{\operatorname{vol}(L_1 \cap T)}{\|\bw_1- \bv\|} \right\rfloor + 1= \left\lfloor \frac{\operatorname{vol}(L_2 \cap T)}{\|\bw_2- \bv\|} \right\rfloor + 1 = |L_2 \cap T \cap \mathbb{Z}^2|.
\]
On the other hand, if $\langle \ba, \bw_1 \rangle < \langle \ba, \bw_2 \rangle$, then there exists a $\bw' \in [\bv, \bw_2]$ such that $\langle \ba, \bw' \rangle = \langle \ba, \bw_1 \rangle$. Again, by proportionality of lines
\[
\frac{\operatorname{vol}(L_1 \cap T)}{\|\bw_1- \bv\|} = \frac{\operatorname{vol}(L_1 \cap T)}{\|\bw'- \bv\|} \geq \frac{\operatorname{vol}(L_2 \cap T)}{\|\bw_2- \bv\|}.
\]
Applying \Cref{lem:nvol} it follows that $|L_1 \cap T \cap \mathbb{Z}^2| \geq |L_2 \cap T \cap \mathbb{Z}^2|$.
% %\input{Figures/2.14_LLgenerator}
\end{proof}

% \begin{remark}
%     Using notation as in \Cref{lemma:min-gens}, if $\langle a, \bw_1 \rangle <\langle a, \bw_2 \rangle$ then we also have strict inequality between the normalized volumes $\nvol(\aff\{\bv,\bw_1\} \cap T) > \nvol(\aff\{\bv,\bw_2\} \cap T)$ however the number of lattice points could be the same.
% \end{remark}

Using the lemmas above, we prove that the algorithm \textsc{Computing Lattice Diameters} is correct and runs in polynomial time in $n$ and $\log_2(c)$, where $n$ is the number of vertices of $P$ and $c$ is the largest absolute value of a coordinate of a vertex.

\begin{theorem}\label{thm:LD-poly}
     Let $P$ be a lattice polygon. There exists a polynomial-time algorithm that computes a lattice diameter line of $P$ for every direction in which a lattice diameter line exists. In particular, the lattice diameter of $P$ can be computed in polynomial time.
\end{theorem}

\begin{proof}
    \emph{Correctness of the algorithm:} Let $\bu$ be a diameter direction of $P$. By \Cref{lemma:vertex_edge_pair}, there exists a lattice diameter line $L$ in direction $\bu$, that passes through an edge $e$ and an opposite vertex $\bv$ to that edge. Thus $L$ is a local lattice diameter line of $T=\conv(\{e,\bv\})$ at $\bv$. Let $m = \min \{3, |T \cap \Z^2 \setminus \{\bv\}|\}$, and let $L_i$ for $i \in [m]$ be the lattice lines computed in the second \textit{for loop}. By Alarcon \cite[Lemma~2.4]{Alarcon_PhD} at most three lattice diameter lines pass through $\bv$, thus $L$ must be contained in $\{L_i \st i \in [m]\}$.

    \emph{Complexity of the algorithm:} We can assume that the input data is a list of vertices, edges, and incidence relations between the vertices and edges, since in dimension two this data can be computed from the vertices of $P$ in polynomial time. Let $n$ denote the number of vertices of $P$.

    \begin{enumerate}
        \item \emph{first for loop:} This loop is repeated $n$ times, since $P$ has $n$ edges. Finding an opposite vertex $\bv$ to an edge $e$ is linear optimization, which can be done in polynomial time in the input \cite{LP_Poly}. Further, checking if $\bv$ lies on an edge $[\bv,\bv']$ that is parallel to $e$, in which case $e$ has two opposite vertices, can also be done in $O(1)$. Each edge contributes at most two triangles, so throughout the whole algorithm $|\Tcal| \leq 2n$.
        \item \emph{second for loop:} This for loop is repeated $|\mathcal{T}|$ times. For one iteration:
        Computing $|T \cap \Z^2|$ can be computed in polynomial time, for example, by Pick's Theorem \cite{pick99}. Let $\ba$ be an outward normal vector of $T$ with $\min_{\bx \in T} \langle \ba, \bx \rangle = \langle \ba,\bv \rangle$.
        A lattice point $\bw_1 \in T \cap \mathbb{Z}^2 \setminus \{\bv\}$ minimizes the functional $\langle \ba, \cdot \rangle$ if and only if it minimizes this functional over the lattice points of the rational polygon
        \[
        T' := T \cap H^-(\ba, \langle \ba, \bv \rangle + 1).
        \]
        The other $\bw_i$ for $i \in [m]$ can be computed similarly, with an additional check if $\bw_i \in T \cap H(\ba, \langle \ba, \bw_{i-1} \rangle)$. Since integer linear programming can be solved in polynomial time in dimension two~\cite{ILP-2d-poly}, the minimizers $\bw_i$ can be computed in polynomial time in the input size of $P$. 
        \item The number of lattice points on a segment $L \cap T$, where $L = \aff(\{\bv,\bw\})$, is given by
        \[
        |L \cap T \cap \mathbb{Z}^2| = \left\lfloor \frac{\operatorname{vol}(L \cap T)}{\|\bw - \bv\|} \right\rfloor + 1.
        \]
        $|\mathcal{L}|$ is bounded above by $3\cdot (2n)$
        % To compute $|L \cap T \cap \Z^2|$: computing the intersection points of $L \cap T$ and then its Euclidean length can be done in constant time. Finding a primitive vector for the direction of the line, reduces to computing greatest common divisors which lies in $O(poly(n\cdot log_2(c)))$. Finally, using \Cref{lem:nvol}, the number of lattice points in $L \cap T$ can also be computed in polynomial-time,
        and finding maximizers of $|L \cap T \cap \Z^2|$ is linear in $|\mathcal{L}|$.
    \end{enumerate}\vspace{-1.5em}
\end{proof}

\begin{remark}
   \Cref{lemma:min-gens} extends to simplices in any dimension. Let $S = \conv(\{F, \bv\})$ be a simplex where $\bv$ is an opposite vertex to a facet $F$. The argument from \mbox{\Cref{thm:LD-poly}} and the fact that integer linear programming can be solved in polynomial time in fixed dimension~\cite{ILP_poly} yield a polynomial-time method to compute the local lattice diameter \mbox{of $S$ at $\bv$}.
\end{remark}

% \begin{remark}
%     Alarcon's Lemma \cite[Lemma~2.4]{Alarcon_PhD} does not say, that a polygon has at most three local lattice diameter lines at any fixed vertex, and this need not be true, as can be seen for $T=\conv\{(0,0),(-m,1), (m,1)\}$ for any $m>1$, where $v=(0,0)$. His lemma says that \emph{if} local lattice diameter lines of $T$ at $v$ are also lattice diameter lines of $T$, \emph{then} there are at most three of them.
% \end{remark}






\subsection{Hardness of computing lattice diameters in dimensions three and higher}\label{sec:NP-hard}

A key ingredient in our algorithm to compute lattice diameters of lattice polygons is the existence of lattice diameter lines that pass through vertices. This fails in higher dimensions, as the following example shows.

\begin{example}\label{ex:LD-3d-interior}
For $\bv_1 = (-1,0), \, \bv_2 = (0,-1), \, \bv_3 = (1,1)$, let
\[
P=\conv (\{ (-m,\bv_1), (-m,\bv_2), (-m+1,\bv_3),(m,-\bv_1),(m,-\bv_2),(m-1,-\bv_3)\}),
\]
where $m \geq 2$. The unique lattice diameter segment of $P$ is the segment on the first-coordinate axis, from values $-(m-1)$ to $m-1$. It does not pass through a vertex of $P$, not even through a lattice point on the boundary of $P$. \Cref{fig:3d-novertex} depicts this polytope for $m=3$ with its lattice diameter segment in orange.

\begin{figure}[ht!]
    \centering
    \tdplotsetmaincoords{80}{15} % {elevation}{rotation around vertcal axis}
\begin{tikzpicture}[tdplot_main_coords, scale=0.9]

% ====== VERTICES ======
\coordinate (A) at (-3,0,-1);  %left-triangle front bottom
\coordinate (B) at (-3,-1,0);  %left-triangle front top
\coordinate (C) at (-2,1,1);   %left-triangle back
\coordinate (D) at (3,0,1);    %-A: right-triangle back top
\coordinate (E) at (3,1,0);    %-B: right-triangle back bottom
\coordinate (F) at (2,-1,-1);  %-C: right-triangle front

%====== BLUE BOXES FOR ORIENTATION ======
% \draw[line width=1mm,blue] (-3,-1,-1)--(-3,-1,1)--(-3,1,1)--(-3,1,-1)--cycle {};
% \draw[line width=1mm,blue] (-2,-1,-1)--(-2,-1,1)--(-2,1,1)--(-2,1,-1)--cycle {};

%outside left x-axis
\draw[line width=0.3mm, black] (-2.67,0,0)--(-3.5,0,0) {}; %x
\fill[black!60] (-3,0,0) circle (2pt);
%z-axis bottom
\draw[line width=0.3mm, black] (0,0,-10/13)--(0,0,-1.5) {};
%\fill[black] (0,0,-1) circle (1.5pt);
%\fill[black] (0,0,-2) circle (1.5pt);

%y-axis back
\draw[line width=0.3mm, black] (0,10/13,0)--(0,6.5,0) {};
% \foreach \yy in {1,...,6} {
%   \fill[black] (0,\yy,0) circle (1.5pt);
% }

% very back face
\fill[lightgray, opacity=0.8] (C)--(D)--(E)--cycle;
% BIGGEST back left face
\fill[lightgray!70, opacity=0.8] (A)--(C)--(E)--cycle;

% bottom right/front face 
\fill[lightgray!30, opacity=0.8] (A)--(E)--(F)--cycle;
% left triangle
\fill[darklav, opacity=0.7] (A)--(B)--(C)--cycle;

%darklav!85!black
\draw[gray!80!black, thick, dashed] (A)--(C)--(E) -- cycle;
\draw[gray!80!black, thick] (A)--(B)--(C);

%z-axis inside
\draw[line width=0.3mm, black] (0,0,-10/13)--(0,0,8/11) {}; 
%y-axis inside
\draw[line width=0.3mm, black] (0,-10/13,0)--(0,8/11,0) {};


% ====== LATTICE DIAMETER ======
%inside part of line
\draw[line width=0.3mm, black] (-2.66,0,0) -- (-2,0,0);
\draw[line width=0.3mm, black] (2,0,0) -- (2.66,0,0);

%convex hull of lattice point 
\draw[very thick, orange, minimum width=4pt] (-2,0,0) -- (2,0,0);
\foreach \xx in {-2,...,2} {
  \fill[orange] (\xx,0,0) circle (2.5pt);
}

% right triangle
\fill[darklav, opacity=0.7] (D)--(E)--(F)--cycle;
\draw[darklav!85!black, thick] (D)--(E)--(F)--cycle;


% ------------ OUTSIDE TOP ------------
% top (left) face
\fill[lightgray!50, opacity=0.2] (B)--(C)--(D)--cycle;
\draw[gray!80!black, thick](B)--(C)--(D)--cycle;

%z-axis top
\draw[line width=0.3mm, black] (0,0,8/11)--(0,0,1.5) {};
%\fill[black] (0,0,1) circle (1.5pt);
%\fill[black] (0,0,2) circle (1.5pt);

% outside right x-axis
\draw[line width=0.3mm, black] (2.67,0,0)--(3.5,0,0) {}; %x
\fill[black] (3,0,0) circle (2pt);

% ------------ OUTSIDE FRONT ------------
% right front face
\fill[lightgray!50, opacity=0.2] (B)--(D)--(F)--cycle;
\draw[gray!90!black, thick](B)--(D)--(F)--cycle;

% very front face
\fill[lightgray!50, opacity=0.2] (A)--(B)--(F)--cycle;
\draw[gray!90!black, thick](A)--(B)--(F)--cycle;

%y-axis front
\draw[line width=0.3mm, black] (0,-10/13,0)--(0,-6.5,0) {};
% \foreach \yy in {-1,...,-6} {
%   \fill[black] (0, \yy,0) circle (1.5pt);
% }


% ====== VERTEX DOTS ======
\foreach \p in {(A),(B),(C),(D), (E), (F)} {
  \node at \p [circle, draw=darklav!85!black, fill=darklav!85!black, minimum width=3pt] {};
}

%====== BLUE BOXES FOR ORIENTATION ======
%\draw[line width=1mm,blue] (4,-1,-1)--(4,-1,1)--(4,1,1)--(4,1,-1)--cycle {};
%\draw[line width=1mm,blue] (3,-1,-1)--(3,-1,1)--(3,1,1)--(3,1,-1)--cycle {};
\end{tikzpicture}
\caption{A \textit{unique} lattice diameter segment contained in the interior of a lattice polytope.}
\label{fig:3d-novertex}
\end{figure}
% \begin{figure}[ht!]
%     \centering
% \includegraphics[width=0.37\linewidth]{Figures/1.3_3d-novertex.png}
%     \caption{A unique lattice diameter segment contained in the interior of $P$.}
%     \label{fig:3d-novertex}
% \end{figure}
\end{example}

We show that in a general setting computing lattice diameters becomes $\NP$-hard.

\begin{theorem}\label{thm:NP-hard}
     Let $d \geq 3$ and let $K \subseteq \R^d$ be a $d$-dimensional bounded semi-algebraic set. \mbox{Computing} the lattice diameter of $K$ is $\NP$-hard.
\end{theorem}

\begin{proof}
For integers $a,b,c>0$ let $f(x,y):=(x^2-a-by)^2$ and $R(a,b,c):=\{ (x,y) \st 1 \leq x \leq c-1, \frac{1-a}{b} \leq y \leq \frac{(c-1)^2-a}{b}\}$.
It was shown in~\cite[Lemma 2.1]{Lemma2.1} that deciding whether the minimum of $f$ over $R \cap \Z^2$ is zero, is $\NP$-hard. This is a polynomial-time reduction from the $\NP$-complete problem AN1 in \cite[p.~249]{GJ}. We may assume that $c>\max\{2,b\}$ (which can be checked in constant time), otherwise $R \cap \Z^2 = \emptyset$. Under these assumptions $R$ is also two-dimensional.

% The problem AN1 is still $\NP$-complete when restricting to these integers since the case $c=1$ has no solutions, and the case $c=2$ asks whether for given $a,b>0$ and $x=1$, we have $x^2=1 \equiv a \pmod b$. 

%If $R \cap \Z^2$ is empty (which can be determined in polynomial-time by linear programming \cite{LP_Poly}) \Anouksays{update to the better 2d reference!}, there is no minimum of $f$ over $R \cap \Z^2$, so .

We construct a three-dimensional semi-algebraic set $K$ with description length polynomial in $a,b,c$, such that determining the lattice diameter of $K$ corresponds to determining the minimum value of $f$ on $R \cap\Z^2$.

\textit{Construction of $K$:} Let $\bp=(1,\lceil \frac{1-a}{b} \rceil ) \in R \cap \Z^2$, $Z:=f(\bp) + \max\{  \lfloor \frac{c^2-2c}{b}\rfloor +1,c-1\}$, and define $ K:= \{ (x,y,z) \st (x,y) \in R, \ f(x,y) \leq z \leq Z \}.$ $K$ is bounded above by $R \times \{Z\}$ and below by the graph of $f$. \Cref{fig:NP-hard-graph} is a sketch of this. 

    \begin{figure}[h!]
\centering
\begin{tikzpicture}[x=0.75pt, y=0.75pt, scale=1]

% ================
% COLORS
% ================

% ========================
% 1. GRID OF LATTICE POINTS
% ========================
% LATTICE POINTS
\foreach \i in {-2,...,13} {           % x-coordinates
  \foreach \j in {-2,...,9} {         % y-coordinates
    \node[circle, fill=gray!70, inner sep=0pt, minimum size=1.5pt] at (20*\i, 20*\j) {};
    % Optionally label:
    % \node[below right, font=\tiny] at (20*\i, 20*\j) {(\i,\j)};
  }
}

% ========================
% 2. COLORED AXES
% ========================
\coordinate (O) at (0,0); % origin now truly at (0,0)

% y
\draw[line width=0.8pt] 
  (O) -- ++(75,50);   
\draw[-{Latex[round]},line width=0.8pt] 
  (O) -- ++(-30,-20);

% x
\draw[-{Latex[round]}, line width=0.8pt] 
  (-35,0) -- ++ (300,0);     

% z
\draw[-{Latex[round]}, line width=0.8pt] 
  (0,-25) -- ++(0,170);    

% ========================
% 3. SHADED POLYGONS & SURFACES
\begin{scope}[yshift=255, xshift=-163]
% Dashed curve
\draw[line width=0.75, dashed] 
  (326.86,-194.09) .. controls (328.18,-203.74) and (329.37,-250.91) ..
  (327.53,-264.4) .. controls (325.68,-277.89) and (323.33,-285.66) ..
  (320.16,-287.83);

%front curve
\draw (308.11,-214.53) .. controls (306.44,-229.86) and (311.56,-275.26) .. (319.56,-287.83) ;

%main body
\draw[draw=black, fill=lightgray, opacity=0.5] 
  (293.44,-207.33) .. controls (294.17,-208.35) and (302.3,-217.08) ..
  (317.82,-213.33) .. controls (333.33,-209.58) and (342.93,-206.33) ..
  (350.93,-205.4) .. controls (358.93,-204.46) and (363.6,-203.13) ..
  (370.4,-206.33) .. controls (377.2,-209.53) and (379.73,-198.86) ..
  (379.6,-199.26) .. controls (379.47,-199.66) and (378.82,-240.46) ..
  (377.02,-256.26) .. controls (375.22,-272.06) and (368.17,-306.66) ..
  (357.83,-308.32) .. controls (347.5,-309.99) and (347.33,-282.13) ..
  (338,-281.63) .. controls (328.67,-281.13) and (322.5,-291.29) ..
  (317.17,-287.29) .. controls (311.83,-283.29) and (310.83,-274.29) ..
  (308.33,-268.29) .. controls (305.83,-262.29) and (292.71,-206.32) ..
  (293.44,-207.33) -- cycle;

  \end{scope}
  
\begin{scope}[yshift=-40, xshift=2]

% Lower lavender panel (wider)
\fill[darklav, opacity=0.25]
   (30,30) -- (150,30) -- (210,70) -- (90,70) -- cycle;

% Upper lavender panel (identical, just shifted up)
\fill[darklav, opacity=0.25]
   (30,30+140) -- (150,30+140) -- 
   (210,70+140) -- (90,70+140) -- cycle;

\end{scope}



% ========================
% 4. ORANGE LATTICE DIAMETER
% ========================
\draw[color=orange, thick] (140,40) -- (140,140);

\foreach \yy in {40,60,80,100,120, 140} {
  \node[circle, fill=orange, inner sep=0pt, minimum size=3pt] at (140,\yy) {};
}

\begin{scope}[yshift=255, xshift=-163]
%top
\draw[black, fill=darklav, fill opacity=0.5] 
  (305.4,-196.07) .. controls (325,-189.72) and (352.23,-198.26) ..
  (368.23,-188.76) .. controls (384.23,-179.26) and (382.82,-212.64) ..
  (369.42,-205.84) .. controls (356.02,-199.04) and (319.71,-217.32) ..
  (306,-213.6) .. controls (292.29,-209.89) and (285.8,-202.41) ..
  (305.4,-196.07) -- cycle ;
\end{scope}

% ========================
% 5. LABELS
% ========================
\node[darklav] at (236,143) {$R \times \{Z\}$};
\node[darklav] at (220,10) {$R$};
\node at (180,90) {$K$};



\end{tikzpicture}
\caption{The construction given in the proof of \Cref{thm:NP-hard}.}
\label{fig:NP-hard-graph}
\end{figure}
% \begin{figure}[ht!]
%     \centering
% \includegraphics[width=0.7\linewidth]{Figures/3.20_NP-hard-graph.png}
%     \caption{$K$ shaded in gray, lattice diameter segment in orange}
%     \label{fig:NP-hard-graph}
% \end{figure}
     
The choice of $Z$ ensures that a small region around the segment $S_\bp:=[(\bp,f(\bp)),(\bp,Z)]$ is contained in the interior of $K$, thus $\dim K =3$.
Furthermore, $S_\bp$ contains at least $\max\{ \lfloor\frac{c^2-2c}{b}\rfloor +1,c-1\} +1$ lattice points. Any segment not in direction $(0,0,1)$ projects to a segment in the $x$-axis, or $y$-axis, but the number of lattice points in the projection of $K$ to these axes is bounded by $c-1$, and $\lfloor \frac{c^2-2c}{b}\rfloor +1$, respectively. Thus a lattice diameter line of $K$ must have direction $(0,0,1)$. Then, for a lattice diameter line $L$ of $K$
\begin{align*}
         |L \cap K \cap \Z^3| = Z-\min_{(x,y) \in R \cap \Z^2}f(x,y) +1.
     \end{align*}
So $|L \cap K \cap \Z^3|=Z-1$ if and only if $\min_{(x,y)\in R \cap \Z^2} f(x,y)=0$, which means that determining the lattice diameter of $K$ is $\NP$-hard.

For $d \geq 3$ consider $K_d:=[0,1]^{d-3} \times K$, which, as a Cartesian product of semi-algebraic sets, is a semi-algebraic set. Analogously, a lattice diameter line of $K_d$ has direction $(0,\ldots,0,1)$, and determining whether the lattice diameter is $Z-1$ is $\NP$-hard, since it determines whether the minimum of $f$ over $R \cap \Z^2$ is zero.

Lastly, note that $K_d$ for $d \geq 3$ is defined by a single quartic polynomial, a constant inequality, which are both polynomial in $a,b,c$, and $2^{d-3}$ linear inequalities in $0$ and $1$ (for $[0,1]^{d-3}$). Hence the description length of $K$ and $K_d$ are polynomial in the encoding length of $a,b$ and $c$.
\end{proof}

\begin{remark}
    The construction of $K$ is such that the lattice diameter is guaranteed to have direction $(0, \ldots, 0,1)$. Thus we proved that computing a lattice diameter of a bounded semi-algebraic set, even in a fixed direction, is $\NP$-hard.
\end{remark}

Furthermore, we conjecture that this computation stays hard, even for lattice polytopes:

\begin{conjecture}
    Let $d \geq 3$ and let $P$ be a lattice $d$-polytope. Computing a lattice diameter of $P$ is an $\NP$-hard problem.
\end{conjecture}


%DETAILS
     % \begin{enumerate} [label=$(\roman*)$]
     %     \item $\dim K =3$, since $f(\bp)<Z$, so the segment between $(\bp,f(\bp))$ and $(\bp,Z)$ is in $K$ and also the graph of a neighborhood of $p$ lies below $R \times \{Z\}$, so their convex hull yields a $3$-dimensional region in $K$.
     %     \item The segment between $(\bp,f(\bp))$ to $(\bp,Z)$ contains at least $\max\{ \lfloor\frac{c^2-2c}{b}\rfloor +1,c-1\} \rfloor +1$ lattice points (the additional $+1$ is because $(\bp,f(\bp))$ is a lattice point itself). Consider a lattice line whose direction is not purely the third coordinate direction. Its lattice points project to lattice points on a line in $R$, whose number is bounded by the side lengths of $R$ plus one, i.e.~bounded by $c-1$ and by $\lfloor \frac{c^2-2c}{b}\rfloor +1$. This means that the direction of any lattice diameter of $K$ must be in direction $\spn\{0,0,1\}$.
     % \end{enumerate}
    
%DETAILS
     % Again any lattice line that is not purely in the third coordinate direction projects to a lattice line in the other coordinate directions which are bounded by the side-lengths of $R$ for the first two coordinates, and by $1$ for the latter coordinates. In all cases the line from $(\bp,f(\bp),0, \ldots ,0)$ to $(\bp,Z,0, \ldots ,0)$ is longer than the projection of the other coordinate directions, which guarantees that a lattice diameter of $K'$ must be purely in the third coordinate direction, i.e.~if $L$ denotes the its length is again
     % \begin{align*}
     %      |L \cap K' \cap  \Z^d| = Z-\min_{(x,y) \in R \cap d\Z^2}f(x,y) +1.
     % \end{align*}
     % So $\min_{(x,y)\in R \cap \Z^2}f(x,y)=0$ if and only if $|L \cap K \cap \Z^d| = Z-1$.


\section{Counting lattice diameter lines}\label{sec:QP}

Ehrhart theory studies the number of lattice points in dilations of a given lattice (or rational) polytope. A fundamental result by Ehrhart states that this counting function agrees with a polynomial function for lattice polytopes, and a quasi-polynomial function for rational polytopes. For details, see \cite{beckrobins}. In this section we show that for lattice polygons $P$ the counting function $\LD_P(k)$, which counts the number of lattice diameter lines (or segments) in the dilate $kP$, agrees with a quasi-polynomial function.

\begin{example}\label{ex:qp}
    Let $P= \conv(\{(0,0), (1,3), (5,1), (6,4)\})$. In \Cref{fig:dilations} we depict the segments $L \cap P$ for lattice diameter lines $L$ of $kP$ for $k=1, \ldots, 6$.
    
    %\input{Figures/2.30_dilations.tex}
\begin{figure}[ht!]
    \centering
    %\includegraphics[width=\linewidth]{Figures/2.30_dilations.png}
    \tikzset{every picture/.style={line width=0.75pt}} % ancho por defecto

\begin{tikzpicture}[x=0.75pt,y=0.75pt,yscale=-.3,xscale=.3]


% ---------- Estilos ----------
\tikzset{
  graydot/.style={circle, draw=lightgray, fill=lightgray, minimum width=1pt, inner sep=0pt},
  sone/.style={circle, draw={rgb,255:red,245;lightgreen,166;blue,35}, fill={rgb,255:red,245;lightgreen,166;blue,35}, minimum width=1pt, inner sep=0pt}, % naranja (S1)
  stwo/.style={circle, draw={rgb,255:red,144;lightgreen,19;blue,254},  fill={rgb,255:red,144;lightgreen,19;blue,254},  minimum width=1pt, inner sep=0pt}  % morado (S2)
}



% ==================================================
% MALLA 1 (izquierda)
% ==================================================
\pgfmathsetmacro{\XL}{460}
\pgfmathsetmacro{\YL}{180}
\foreach \i in {0,...,6}{
  \foreach \j in {0,...,4}{
    \node[graydot] at ({\XL + 20*\i},{\YL + 20*\j}) {};
  }
}
% Quadrilateros
\foreach \j in {0}{
  \pgfmathsetmacro{\XL}{460+\j*20}
  \pgfmathsetmacro{\YL}{180-\j*60}
  \foreach \i in {0}{
    \draw [color=black, draw opacity=1]
  (\XL+\i*100,\YL+80-\i*20)  -- (\XL+100+\i*100,\YL+60-\i*20) -- (\XL+120+\i*100,\YL-\i*20) --(\XL+20+\i*100,\YL+20-\i*20) -- cycle;
  }
}

% Dibujar líneas naranjas paralelas
\draw[color=orange, line width=.8pt, draw opacity=1] (478, 200) -- (576, 200);
\draw[color=orange, line width=.8pt, draw opacity=1] (471, 220) -- (569, 220);
\draw[color=orange, line width=.8pt, draw opacity=1] (464, 240) -- (562, 240);

  

% ==================================================
% MALLA 2 (medio)  ---- puntos grises con bucle
% ==================================================
\pgfmathsetmacro{\XL}{1020}
\pgfmathsetmacro{\YL}{140}
\foreach \i in {0,...,12}{
  \foreach \j in {0,...,8}{
    \node[graydot] at ({\XL + 20*\i},{\YL + 20*\j}) {};
  }
}

% Quadrilateros
\foreach \j in {0,1}{
  \pgfmathsetmacro{\XL}{1020+\j*20}
  \pgfmathsetmacro{\YL}{220-\j*60}
  \foreach \i in {0,1}{
    \draw [color=black, draw opacity=1]
  (\XL+\i*100,\YL+80-\i*20)  -- (\XL+100+\i*100,\YL+60-\i*20) -- (\XL+120+\i*100,\YL-\i*20) --(\XL+20+\i*100,\YL+20-\i*20) -- cycle;
  }
}

% Dibujar líneas moradas paralelas
\draw[color=darklav, line width=.8pt, draw opacity=1] (1056, 180) -- (1247.61, 180);
\draw[color=darklav, line width=.8pt, draw opacity=1] (1051, 200) -- (1242.61, 200);
\draw[color=darklav, line width=.8pt, draw opacity=1] (1037, 240) -- (1229.61, 240);
\draw[color=darklav, line width=.8pt, draw opacity=1] (1030, 260) -- (1222.61, 260);

% ==================================================
% MALLA 3 (derecha)  ---- puntos grises con bucle
% ==================================================
\pgfmathsetmacro{\Xo}{1680}
\pgfmathsetmacro{\Yo}{110}
\foreach \i in {0,...,18}{
  \foreach \j in {0,...,12}{
    \node[graydot] at ({\Xo + 20*\i},{\Yo + 20*\j}) {};
  }
}

% Quadrilateros
\foreach \j in {0,1,2}{
  \pgfmathsetmacro{\XL}{1680+\j*20}
  \pgfmathsetmacro{\YL}{270-\j*60}
  \foreach \i in {0,1,2}{
    \draw [color=black, draw opacity=1]
  (\XL+\i*100,\YL+80-\i*20)  -- (\XL+100+\i*100,\YL+60-\i*20) -- (\XL+120+\i*100,\YL-\i*20) --(\XL+20+\i*100,\YL+20-\i*20) -- cycle;
  }
}

\draw[color=lightgreen, line width=.8pt, draw opacity=1] (1740, 170) -- (2020, 170);
\draw[color=lightgreen, line width=.8pt, draw opacity=1] (1720, 230) -- (2000, 230);
\draw[color=lightgreen, line width=.8pt, draw opacity=1] (1700, 290) -- (1980, 290);


% ==================================================
% MALLA 4 (izquierda) ---- puntos grises con bucle
% ==================================================
\pgfmathsetmacro{\XL}{290}
\pgfmathsetmacro{\YL}{600}
\foreach \i in {0,...,24}{
  \foreach \j in {0,...,16}{
    \node[graydot] at ({\XL + 20*\i},{\YL + 20*\j}) {};
  }
}

% Quadrilateros
\foreach \j in {0,...,3}{
  \pgfmathsetmacro{\XL}{290+\j*20}
  \pgfmathsetmacro{\YL}{840-\j*60}
  \foreach \i in {0,...,3}{
    \draw [color=black, draw opacity=1]
  (\XL+\i*100,\YL+80-\i*20)  -- (\XL+100+\i*100,\YL+60-\i*20) -- (\XL+120+\i*100,\YL-\i*20) --(\XL+20+\i*100,\YL+20-\i*20) -- cycle;
  }
}

% Dibujar líneas naranjas paralelas
\draw[color=orange, line width=.8pt, draw opacity=1] (371, 680) -- (745.61, 680);
\draw[color=orange, line width=.8pt, draw opacity=1] (364, 700) -- (738.61, 700);
\draw[color=orange, line width=.8pt, draw opacity=1] (357, 720) -- (731.61, 720);
\draw[color=orange, line width=.8pt, draw opacity=1] (350, 740) -- (724.61, 740);
\draw[color=orange, line width=.8pt, draw opacity=1] (343, 760) -- (717.61, 760);
\draw[color=orange, line width=.8pt, draw opacity=1] (336, 780) -- (710.61, 780);
\draw[color=orange, line width=.8pt, draw opacity=1] (329, 800) -- (703.61, 800);
\draw[color=orange, line width=.8pt, draw opacity=1] (322, 820) -- (696.61, 820);
\draw[color=orange, line width=.8pt, draw opacity=1] (315, 840) -- (689.61, 840);


% ==================================================
% MALLA 5 (medio)  ---- puntos grises con bucle
% ==================================================
\pgfmathsetmacro{\Xo}{850}
    \pgfmathsetmacro{\Yo}{560}
\foreach \i in {0,...,30}{
  \foreach \j in {0,...,20}{
    \node[graydot] at ({\Xo + 20*\i},{\Yo + 20*\j}) {};
  }
}


% Quadrilateros
\foreach \j in {0,...,4}{
  \pgfmathsetmacro{\XL}{850+\j*20}
  \pgfmathsetmacro{\YL}{880-\j*60}
  \foreach \i in {0,...,4}{
    \draw [color=black, draw opacity=1]
  (\XL+\i*100,\YL+80-\i*20)  -- (\XL+100+\i*100,\YL+60-\i*20) -- (\XL+120+\i*100,\YL-\i*20) --(\XL+20+\i*100,\YL+20-\i*20) -- cycle;
  }
}

% Dibujar líneas moradas paralelas
\draw[color=darklav, line width=.8pt, draw opacity=1] (948, 660) -- (1417.61, 660);
\draw[color=darklav, line width=.8pt, draw opacity=1] (943, 680) -- (1412.61, 680);
\draw[color=darklav, line width=.8pt, draw opacity=1] (929, 720) -- (1399.61, 720);
\draw[color=darklav, line width=.8pt, draw opacity=1] (922, 740) -- (1392.61, 740);
\draw[color=darklav, line width=.8pt, draw opacity=1] (908, 780) -- (1378.61, 780);
\draw[color=darklav, line width=.8pt, draw opacity=1] (901, 800) -- (1371.61, 800);
\draw[color=darklav, line width=.8pt, draw opacity=1] (888, 840) -- (1357.61, 840);
\draw[color=darklav, line width=.8pt, draw opacity=1] (883, 860) -- (1353.61, 860);


% ==================================================
% MALLA 6 (derecha)  ---- puntos grises con bucle
% ==================================================
\pgfmathsetmacro{\Xo}{1540}
\pgfmathsetmacro{\Yo}{520}
\foreach \i in {0,...,36}{
  \foreach \j in {0,...,24}{
    \node[graydot] at ({\Xo + 20*\i},{\Yo + 20*\j}) {};
  }
}

% Quadrilateros
\foreach \j in {0,...,5}{
  \pgfmathsetmacro{\XL}{1540+\j*20}
  \pgfmathsetmacro{\YL}{920-\j*60}
  \foreach \i in {0,...,5}{
    \draw [color=black, draw opacity=1]
  (\XL+\i*100,\YL+80-\i*20)  -- (\XL+100+\i*100,\YL+60-\i*20) -- (\XL+120+\i*100,\YL-\i*20) --(\XL+20+\i*100,\YL+20-\i*20) -- cycle;
  }
}

% Dibujar líneas moradas paralelas

\draw[color=lightgreen, line width=.8pt, draw opacity=1] (1662, 640) -- (2224, 640);
\draw[color=lightgreen, line width=.8pt, draw opacity=1] (1642, 700) -- (2204, 700);
\draw[color=lightgreen, line width=.8pt, draw opacity=1] (1622, 760) -- (2184, 760);
\draw[color=lightgreen, line width=.8pt, draw opacity=1] (1600, 820) -- (2164, 820);
\draw[color=lightgreen, line width=.8pt, draw opacity=1] (1582, 880) -- (2144, 880);








% Text Node
\draw (490,300) node [anchor=north west][inner sep=0.75pt]    {$P$};
% Text Node
\draw (1120,335) node [anchor=north west][inner sep=0.75pt]    {$2P$};
% Text Node
\draw (1840,385) node [anchor=north west][inner sep=0.75pt]    {$3P$};
% Text Node
\draw (480,960) node [anchor=north west][inner sep=0.75pt]    {$4P$};
% Text Node
\draw (1120,1010) node [anchor=north west][inner sep=0.75pt]    {$5P$};
% Text Node
\draw (1840,1050) node [anchor=north west][inner sep=0.75pt]    {$6P$};



\end{tikzpicture}

    \caption{The pattern of lattice diameter lines repeats periodically under dilation of $P$.}
    \label{fig:dilations}
\end{figure}

    Then for all $k \in \N$ we get\vspace{-1em}
    \begin{align*}
        \LD_P(k)=\begin{cases}
            2k+1 & \text{if } k \equiv 1 \pmod 3,\\
            \frac{4}{3} k + \frac{4}{3} & \text{if } k \equiv 2 \pmod 3,\\
            \frac{2}{3} k + 1 & \text{if } k \equiv 0 \pmod 3.
        \end{cases}
    \end{align*}
    Indeed, $\LD_P(1)=3,\ \LD_P(2)=4,\ \LD_P(3)=3,\ \LD_P(4)=9,\ \LD_P(5)=8$ and $\LD_P(6)=5$.
\end{example}

For a set $X \subseteq \R^d$ we write $kX:= \{k \bx \st \bx \in X\}$ for the dilated set by $k \in \N$. In particular, for a polytope $P$, any line intersecting $kP$ can be written as the dilated line $kL$ for a \mbox{line $L$ intersecting $P$.}

\begin{lemma}\label{lemma:same-dir}
    Let $P$ be a lattice polygon and let $k \in \N$. Then, a diameter direction of $kP$ is also a diameter direction of $P$.
\end{lemma}
\begin{proof}
    Let $\bu$ be a diameter direction of $kP$ and let $kL$ be a lattice diameter line with direction $\bu$ that passes through a vertex of $kP$. (That is, $L$ is a line that passes through the corresponding vertex of $P$). Let $L'$ be a lattice line that passes through a vertex of $P$. Then, $kL'$ also passes through a vertex of $kP$ and by \Cref{lem:nvol}
    \begin{align}\label{eq:diam-dir-stay}
        \lfloor \nvol(kL \cap kP) \rfloor +1 = |kL \cap kP \cap \Z^2|  \geq |kL' \cap kP \cap \Z^2| = \lfloor \nvol(kL' \cap kP) \rfloor +1.
    \end{align}
    The normalized volume is a homogeneous function in $k$, so \eqref{eq:diam-dir-stay} and \Cref{lem:nvol} yield
    %since $L,L'$ pass through vertices of $P$ we get
    \begin{align*}
        |L \cap P \cap \Z^2|  = \lfloor \nvol(L \cap P) \rfloor +1 \geq \lfloor \nvol(L' \cap P) \rfloor +1 =|L' \cap P \cap \Z^2|.
    \end{align*}
    Since for every diameter direction of $P$ there exists a lattice diameter line in that direction, that passes through a vertex, we proved that $L$ is a lattice diameter line of $P$. In particular, $\bu$ is a diameter direction of $P$.
\end{proof}

% The next lemma says that only some lattice diameter segments ``remain'' a lattice diameter segment under dilation of the polygon, precisely those for which $L \cap P$ has largest normalized length. \Cref{fig:nvol-dom} illustrates this phenomenon.
The following lemma shows that under dilation of the polygon only certain lattice diameter lines remain, specifically those for which $L \cap P$ has the largest normalized length. \Cref{fig:nvol-dom} illustrates this phenomenon.

%dilated polytope 3x where ldsegments disappear

\begin{figure}[ht!]
    \centering
    \begin{tikzpicture}[thin, scale=0.28]\vspace{0.5cm}
        \foreach \x in {-4,...,36} {
            \foreach \y in {-1,...,13} {
                \node at (\x,\y) [circle, draw=lightgray, fill= lightgray, minimum width=1.5pt] {};   
            }
        }
    
    %Outline of polytopes 1P, 2P, 3P
    \draw[line width=0.3mm, black] (-2,0) -- (-3,2) --(-1,4) --(1,4) -- (2,1)-- (0,0) -- cycle;
    \draw[line width=0.3mm, black](8,0) -- (6,4) -- (10,8) -- (14,8) -- (16,2) -- (12,0) -- cycle;
    \draw[line width=0.3mm, black] (23,0) -- (20,6) -- (26,12) -- (32,12) -- (35,3) -- (29,0) -- cycle;

    
    %Lattice Diameters in P
    %full segments L \cap P
    % \draw[line width=0.3mm, orange] (-3,2) -- (1.66,2);
    % \draw[line width=0.3mm, orange] (-2.5,1) -- (2,1);
    
    %ends at last lattice point
    \draw[line width=0.3mm, orange] (-3,2) -- (1,2);
    \draw[line width=0.3mm, orange] (-2,1) -- (2,1);
    
    %same for both
    \draw[line width=0.3mm, orange] (-1,0) -- (-1,4);
    \draw[line width=0.3mm, orange] (0,0) -- (0,4);

     \foreach \x in {-2, ..., 2} {
    \node at (\x,1) [circle, fill=orange, minimum width=3pt] {};
        }
    \foreach \x in {-3, ..., 1} {
    \node at (\x,2) [circle, fill=orange, minimum width=3pt] {};
        }
    \foreach \y in {0, ..., 4} {
    \node at (-1,\y) [circle, fill=orange, minimum width=3pt] {};
    \node at (0,\y) [circle, fill=orange, minimum width=3pt] {};
        }

        
    %Lattice Diameters in 2P
    %\draw[line width=0.3mm, orange] (6,4) -- (15.32,4);
    \draw[line width=0.3mm, orange] (6,4) -- (15,4);
    \draw[line width=0.3mm, orange] (7,2) -- (16,2);

    \foreach \x in {7,8, ..., 16} {
        \node at (\x,2) [circle, fill=orange, minimum width=3pt] {};
    }
    \foreach \x in {6,7,..., 15} {
        \node at (\x,4) [circle, fill=orange, minimum width=3pt] {};
    }


    %Lattice Diameters in 3P  
    \draw[line width=0.3mm, orange] (20,6) -- (34,6);

    \foreach \x in {20, ..., 34} {
        \node at (\x,6) [circle, fill=orange, minimum width=3pt] {};
    }
    
    \end{tikzpicture}
    \caption{Lattice diameter segments can ``disappear'' when dilating a polygon.}
    \label{fig:nvol-dom}
\end{figure}
% \begin{figure}[ht!]
%     \centering
%     \includegraphics[width=0.65\linewidth]{Figures/2.30_nvol-dom.png}
%     \caption{lattice diameter segments when dilating a polygon}
%     \label{fig:nvol-dom}
% \end{figure}

\begin{lemma}\label{lemma:same-nvol}
    Let $P$ be a lattice polygon. There exists a positive integer $q$ such that for all $k\geq q$, and all lattice diameter lines $kL,kL'$ of $kP$, $\nvol(L \cap P)=\nvol(L' \cap P)$.
\end{lemma}

\begin{proof}
    Let $\Rcal$ be the set of rational lines whose direction is a diameter direction of $P$.
    By \Cref{lemma:vertex_edge_pair} we can choose a lattice diameter line $L \in \Rcal$ such that $\nvol(L \cap P) \geq \nvol(L' \cap P)$ for all rational lines $L' \in \Rcal$. Let $\nvol(L \cap P)=\frac{p}{q}$ for some $p,q \in \N$ and note that $\frac{p}{q} \geq 1$.

    Let $kL'$ be a lattice diameter line of $kP$. By \Cref{lemma:same-dir} its direction is a diameter direction of $P$, so $L' \in \Rcal$. It suffices to show $\nvol(L \cap P) \leq \nvol(L' \cap P)$, so that we obtain equality. Towards a contradiction, suppose  $\nvol(L \cap P) > \nvol(L' \cap P)$. The normalized volume is a homogeneous function in $k$, so 
    \begin{align*}
        \nvol(kL' \cap kP) < \nvol(kL \cap kP) = p + \frac{(k-q)p}{q}.
    \end{align*}
    Together with \Cref{lem:nvol} and using $\frac{(k-q)p}{q}\geq 1$, this yields
    \begin{align*}
        |kL' \cap kP \cap \Z^2| \leq \lfloor \nvol(kL' \cap kP)\rfloor +1 < p  + \left\lfloor \frac{(k-q)p}{q}\right\rfloor +1 = |kL \cap kP \cap \Z^2|,
    \end{align*}
    a contradiction to $kL'$ being a lattice diameter line of $kP$.
\end{proof}

Next we observe that parallel lattice diameter lines, for which the segment $\vol(L \cap P)$ is maximal, occur between two parallel edges of $P$.

\begin{lemma}\label{lemma:LD-in-chambers}
    Let $P$ be a lattice polygon and let $\bu \in \Z^2 \setminus \{0\}$. Define
    \[
    \Lcal = \{ L \st \text{ L is a $\bu$-lattice line and } \vol(L \cap P) \geq \vol(L' \cap P) \text{ for all $\bu$-lattice lines $L'$}\}.
    \]
    If $|\Lcal|>1$, then there exist parallel edges $e_1,e_2$ of $P$ such that $L \subseteq \conv(e_ 1 \cup e_2)$ for all $L \in \Lcal$.
\end{lemma}

\begin{proof}
    Let $\ba \in \Z^2\setminus \{0\}$ be orthogonal to $\bu$ and let $\bv_1, \ldots, \bv_n$ be the vertices of $P$, enumerated such that $\langle \ba, \bv_i \rangle \leq \langle \ba,\bv_{i+1} \rangle$ for all $i \in [n-1]$. Define the \textit{parallel chambers}
    \[
    C_i:= P \cap H^+(\ba,\langle \ba,\bv_i \rangle ) \cap H^-(\ba,\langle \ba,\bv_{i+1} \rangle )
    \]
    for all $i \in [n-1]$ for which $\langle \ba,\bv_i \rangle < \langle \ba , \bv_{i+1} \rangle$, so $C_i$ is full-dimensional. A chamber $C_i$ is bounded by $H(\ba,\langle \ba, \bv_i\rangle), \, H(\ba,\langle \ba, \bv_{i+1}\rangle)$, and two edges $e_{i1},e_{i2}$ of $P$. If the edges $e_{i1}$ and $e_{i2}$ are parallel, then all segments $C_i \cap H(\ba,b_i)$ for $b_i \in [ \langle \ba,\bv_i \rangle, \langle \ba,\bv_{i+1} \rangle]$ have the same length. Otherwise, there is a unique segment of the form $C_i \cap H(\ba,b_i)$ where $b_i \in [ \langle \ba,\bv_i \rangle, \langle \ba,\bv_{i+1} \rangle]$, which is the longest among all segments in this chamber. Note that $P \cap H(\ba,b_i) = C_i \cap H(\ba,b_i)$, so if in the setting of this lemma $|\Lcal|>1$, then all $L \cap P$ for $L \in \Lcal$ have the same length, and thus are contained in a chamber bounded by parallel edges.
\end{proof}

Finally, we count the number of $\bu$-lattice diameter lines in dilations of a rational parallelogram. For $\bu \in \Z^{2} \setminus \{0\}$, define $LD_{P,u}(k)$ as the number of $\bu$-lattice diameter lines (or segments) of $kP$.

\begin{proposition}\label{prop:QP-in-parallelotope}
    Let $R$ be a rational parallelogram with parallel edges in direction $\bu$ that each contain an integral vertex. Then there exists a positive integer $q$ such that \mbox{$\LD_{R,u}:\N_{\geq q} \to \N$} agrees with a quasi-polynomial function of degree one whose period divides $q$.
\end{proposition}
\begin{proof}
    After applying a unimodular transformation, we may assume that $(0,0)$ is a vertex of $R$ and that $\bu = (1,0)$. Let the other two parallel edges, adjacent to those in direction $\bu$, have direction $(a,q)$ for some $a,q \in \Z, q >0$. Then, one of the edges in direction $\bu$ has the form $[(0,0),(\frac{p}{q},0)]$ and without loss of generality $p>0$. This edge contains a $\bu$-lattice diameter segment. Furthermore, the other edge emanating from the origin has the form $[(0,0),\lambda(a,q)]$ for some $\lambda = \frac{w}{q}$ and again we can assume $w \in \N$. That is, $\lambda(a,q) = (\frac{aw}{q},w)$.
    
    Let $k \in \N$ be a dilation factor. The lattice width of $kR$ in direction $(0,1)$ is $kw$.
    %The normalized length of lattice segments in direction $(0,1)$ agrees with their Euclidean length, since $\|(0,1)\|=1$, but we continue to write $\nvol$, because the formulas hold in general for the normalized length.
    For $j \in \{0, \ldots, kw\}$, define the lines on height $j$ as $L_j = (0,j) + \spn(\{(1,0)\})$ and denote the corresponding segments by $S_j(k) = L_j \cap kR$. When $k=1$ we just write $S_j$. Let $\Lcal(k) = \{S_j(k) \st j \in \{0, \ldots, kw\}\}$, that is, all $\bu$-lattice segments that intersect $kR$. For a segment $S \subseteq \R^2$, let $S_{{\Z}} := \conv (\{S \cap \Z^2\})$. So $S_j(k)$ contains a $\bu$-lattice diameter segment if and only if $S_j(k)_{\Z}$ is a $\bu$-lattice diameter segment.

    \emph{Claim 1:} $\nvol(S_j(k))-\nvol(S_{j}(k)_{\Z}) <1$ if and only if $S_{j}(k)_\Z$ is a $\bu$-lattice diameter segment.
    
    It suffices to show this for $k=1$ since $R$ is arbitrary. All $S_j$ have the same length, and $S_0$ is a lattice diameter segment since it contains a vertex. By \Cref{lem:nvol}
    \begin{align*}
        |S_0 \cap \Z^2| = \lfloor \nvol(S_0) \rfloor +1 
        = \lfloor \nvol(S_j) \rfloor +1  \geq |S_j \cap \Z^2| = \nvol({S_{j}}_{\Z}) +1.
    \end{align*}
    Thus $|S_0 \cap \Z^2| = |S_j \cap \Z^2|$ if and only if $\lfloor \nvol(S_j) \rfloor +1  = \nvol({S_{j}}_{\Z}) +1$ which holds if and only if $\nvol(S_j)- \nvol(S_{j{\Z}}) < 1$. This proves \textit{Claim 1}.


    \textit{Claim 2:} $\bu$-lattice diameter segments $S_j(k)_{\Z}$ are determined by $j \pmod q$. 
    
    Let $\chi_k: \{0, \ldots, kw\} \to \{0,1\}$ be the indicator function with $\chi_k(j)=1$ if and only if $L_j$ is a $\bu$-lattice diameter line of $kR$. Let $j,j' \in \{0, \ldots, kw\}$ with $j'= j+tq$ for some $t \in \N$. Then $S_{j'}(k) =S_j(k)+(ta,tq)$, and $S_{j'}(k)_\Z =S_j(k)_{\Z}+(ta,tq)$. Thus
    \begin{align*}
         \nvol(S_j(k))-\nvol(S_{j}(k)_{\Z})=\nvol(S_{j'}(k))-\nvol(S_{j'}(k)_{\Z})
    \end{align*}
    and by \textit{Claim 1} $S_j(k)_\Z$ is a $\bu$-lattice diameter segment of $kR$ if and only if $S_{j'}(k)_\Z$ is.
    
    \textit{Claim 3:} $\bu$-lattice diameter segments $S_j(k)_\Z$ are determined by $k \pmod q$.
    
    Let $j \in \{0, \ldots, kw\}$ and $i \in \{0, \ldots, q-1\}$. If $S_j(k) = [\ba,\bb]$, then $S_j(i) = [\ba,\bb + ((k-i)\frac{p}{q},0)]$. When $k \equiv i \pmod q$ the right endpoints of $S_j(k)$ and $S_j(i)$ differ by a vector in $\Z \times \{0\}$, so
    \begin{align*}
         \nvol(S_j(k))-\nvol(S_j(k)_\Z)=\nvol(S_j(i))-\nvol(S_j(i)_\Z).
    \end{align*}
    By \textit{Claim 1} it follows that $S_j(k)_\Z$ is a $\bu$-lattice diameter segment of $kR$ if and only if $S_j(i)_\Z$ is a $\bu$-lattice diameter segment of $iR$.
    
    \textit{The counting function:} Note that $q$ and $w$ are values that only depend on $R$. Let $i \in \{0, \ldots, q-1\}$, and let $k=\tilde{k}q+i$ for some $\tilde{k} \in \N$. Define a $q$-block of $kR$ as the set of lines $L_{tq}, \ldots, L_{tq+q-1}$ for those $t \in \N$ for which $tq+q-1 \leq kw$. The number of $q$-blocks in $kR$ is
    \begin{align*}
         \blocks_i(k)  = \left \lfloor  \frac{kw+1}{q} \right\rfloor = \tilde{k}w + \left\lfloor \frac{iw+1}{q}\right\rfloor = \frac{wk}{q}  - \lp \frac{iw}{q} - \left\lfloor \frac{iw+1}{q} \right\rfloor \rp,
    \end{align*}
    which is a linear function in $k$. Next we count the number of $\bu$-lattice diameter lines in $kR$ per $q$-block. By \textit{Claim 3}, this only depends on the congruence class $k \bmod q$. Take $q+i$ as a representative for $i \bmod q$,
    % The width of $(q+i)R$ in direction $(0,1)$ is $(q+i)w>q$, so $(q+i)R$ 
    since it has at least one $q$-block, and let
    \begin{align*}
        n_i := \sum_{j=0}^{q-1}\chi_{q+i}(j) = \# \{  L_j: j =0, \ldots, q-1,\, L_j \text{ is a $\bu$-lattice diameter line of } (q+i)P\}.
    \end{align*}
    Define $\rem_i$ as the number of lattice lines intersecting $kR$ not contained in a $q$-block, then $\rem_i= iw+1 \bmod q$. Clearly $kw +1 \equiv iw+1 \pmod q$. 
    The number of those lines that are $\bu$-lattice diameter lines is
    \begin{align*}
        r_i := \sum_{j=0}^{\rem_i} \chi_{i}(j)  = \# \{  L_j: j =0, \ldots, \rem_i, \, L_j \text{ is a $\bu$-lattice diameter line of } iP\}.
    \end{align*}
    Together we get $\LD_{R,u\,}(k) = n_i \cdot \blocks_{i}(k) + r_i$ which for fixed congruence class $ i \bmod q$ is a linear function in $k$. Thus $\LD_{R,u}$ agrees with a quasi-polynomial function of degree one, whose period \mbox{divides $q$}.
\end{proof}

Now we are ready to prove the main theorem of this subsection.

\begin{theorem}\label{thm:QP}
    Let $P$ be a lattice polygon. There exists a positive integer $q$ such that \mbox{$\LD_P: \N_{\geq q} \to \N$} agrees with a quasi-polynomial function of degree one whose period divides $q$.
\end{theorem}


\begin{proof}
    Let $L$ be a lattice diameter line of $P$ such that $\nvol(L \cap P)$ is maximal among all lattice diameter lines of $P$. Let $\nvol(L \cap P) =\frac{p}{q}$ for $p,q \in \N$. By \Cref{lemma:same-dir} there exists a set $\Dcal$ that is the set of diameter directions of $kP$ for all $k\geq q$.
    By \cite[Theorem~4.1]{Alarcon_PhD}, $|\Dcal| \leq 4$.

    Let $k \geq q$. By \Cref{lemma:same-nvol}, $\nvol(L \cap P)$ is the same for all lattice diameter lines $L$ of $kP$. Let $\Dcal = \Dcal_{=1} \cup \Dcal_{>1}$, where $\bu \in \Dcal_{=1}$ if $kP$ has exactly one $\bu$-lattice diameter line and otherwise $\bu \in \Dcal_{>1}$. As in the proof of \Cref{lemma:LD-in-chambers}, for $\bu \in \Dcal_{>1}$, there exist vertices $\bv_i, \bv_{i+1}$ of $P$ such that all $\bu$-lattice diameter lines are contained in the dilated rectangle $k\cdot R(\bu)$ for $R(\bu):=P \cap H^+(\ba,\langle \ba, \bv_i \rangle) \cap H^-(\ba,\langle \ba, \bv_{i+1} \rangle)$, where $\ba \in \bu^{\perp} \setminus \{0\}$. Then
    \begin{align*}
        \LD_P(k) = \sum_{\substack{R = R(\bu) \\ \bu \in \Dcal_{>1}}} \LD_{R,\bu\,}(k) + \sum_{\substack{\bu \in \Dcal_{=1}}} 1
    \end{align*}
    and by \Cref{prop:QP-in-parallelotope}, $\LD_{R,\bu\,}(k)$ agrees with a quasi-polynomial function of degree one whose period divides $q$, so $\LD_P(k)$ does as well.
\end{proof}

We return to \Cref{ex:qp} and demonstrate how we obtain the claimed quasi-polynomial.

\begin{example}\label{ex:qp-with-proof}
    $P = \conv (\{ (0,0),(1,3), (5,1), (6,4)\})$ and its dilates $2P$ and $3P$ are depicted in \Cref{fig:dilations-123} with their lattice diameter segments. $P$ has one diameter direction $\bu=(1,0)$, so \mbox{$\Dcal = \{(1,0)\}$}. All lattice diameter segments of $P$ are contained in the parallelogram $R = \conv (\{ (1/3,1), (1,3), (17/3,3)\})$. The slope of the left edge is $3$ which implies that the period of $\LD_{R,\bu}$ is $q=3$, and the lattice width of $R$ in direction $(0,1)$ is $w=2$. \Cref{thm:QP} guarantees a quasi-polynomial behavior for values $k \geq q$, but here we actually get it for all $k \in \N$.

    \begin{figure}[ht!]
\centering
% --------- P ---------
\begin{minipage}[b]{0.33\textwidth}
\raggedright
\begin{tikzpicture}[scale=0.6]\vspace{0.5cm}
\hspace{1em}
\foreach \x in {0,...,6} {
    \foreach \y in {0,...,4} {
        \node at (\x,\y) [fill=lightgray, minimum width=2pt] {};   
    }
}
\draw[line width=0.3mm, black] (0,0) -- (1,3) -- (6,4) -- (5,1) -- cycle;

\draw[line width=0.3mm, orange] (1,1) -- (5,1);
\draw[line width=0.3mm, orange] (1,2) -- (5,2);
\draw[line width=0.3mm, orange] (1,3) -- (5,3);


\foreach \x in {1,...,5} {
    \node at (\x,1) [circle, fill=orange, minimum width=3pt] {};
    \node at (\x,2) [circle, fill=orange, minimum width=3pt] {};
    \node at (\x,3) [circle, fill=orange, minimum width=3pt] {};
}
\end{tikzpicture}
\end{minipage}%
% --------- 2P ---------
\begin{minipage}[b]{0.33\textwidth}
\centering
\begin{tikzpicture}[scale=0.3]\vspace{0.5cm}
\foreach \x in {0,...,12} {
    \foreach \y in {0,...,8} {
        \node at (\x,\y) [circle, draw=white, fill=lightgray, minimum width=2.5pt] {};   
    }
}
\draw[line width=0.3mm, black] (0,0) -- (2,6) -- (12,8) -- (10,2) -- cycle;


\draw[line width=0.3mm, darklav] (1,2) -- (10,2);
\draw[line width=0.3mm, darklav] (1,3) -- (10,3);
\draw[line width=0.3mm, darklav] (2,5) -- (11,5);
\draw[line width=0.3mm, darklav] (2,6) -- (11,6);

\foreach \x in {1,...,10} {
    \node at (\x,2) [circle, fill=darklav, minimum width=3pt] {};
    \node at (\x,3) [circle, fill=darklav, minimum width=3pt] {};
}
\foreach \x in {2,...,11} {
    \node at (\x,5) [circle, fill=darklav, minimum width=3pt] {};
    \node at (\x,6) [circle, fill=darklav, minimum width=3pt] {};
}

\end{tikzpicture}
\end{minipage}%
% --------- 3P ---------
\begin{minipage}[b]{0.33\textwidth}
\raggedleft
\begin{tikzpicture}[scale=0.2]\vspace{0.5cm}
\foreach \x in {0,...,18} {
    \foreach \y in {0,...,12} {
        \node at (\x,\y) [circle, draw=white, fill=lightgray, minimum width=2.2pt] {};   
    }
}
\draw[line width=0.3mm, black] (0,0) -- (3,9) -- (18,12) -- (15,3) -- cycle;


\draw[line width=0.3mm, lightgreen] (1,3) -- (15,3);
\draw[line width=0.3mm, lightgreen] (2,6) -- (16,6);
\draw[line width=0.3mm, lightgreen] (3,9) -- (17,9);

\foreach \x in {1,...,15} {
    \node at (\x,3) [circle, fill=lightgreen, minimum width=3pt] {};
}
\foreach \x in {2,...,16} {
    \node at (\x,6) [circle, fill=lightgreen, minimum width=3pt] {};
}
\foreach \x in {3,...,17} {
    \node at (\x,9) [circle, fill=lightgreen, minimum width=3pt] {};
}
\end{tikzpicture}
\hspace{1em}
\end{minipage}
\caption{Lattice diameter segments in dilates of $P$ to visualize the values $n_i,\rem_i,r_i$ and $\blocks_i$.}
\label{fig:dilations-123}
\end{figure}
    %  \begin{figure}[ht!]
    %     \centering
    % \includegraphics[width=0.9\linewidth]{Figures/2.30_dilations.png}
    %     \caption*{}
    %     \label{fig:dilations-123}
    % \end{figure}
    
    For $i \in \{0, 1,2\}$ and $k \equiv i \pmod 3$ we partition the $\bu$-lattice lines that intersect $kR$ into $\blocks_i(k)$ $q$-blocks and $\rem_i$ remaining lattice lines. Here, $R$ and $2R$ have one $q$-block, and $3R$ has two $q$-blocks. These examples suffice to read off $n_i$, the number of lattice diameter lines per $q$-block, and $r_i$, the number of lattice diameter lines in the remaining $\rem_i$ lines, of $kR$ for any $k$. We get
  \begin{table}[h!]
    \centering
    \begin{tabular}{c|c|c|c|c}
    $i$ & $n_i$ & $\blocks_i(k)$ &$\rem_i$ & $r_i$\\ \hline \rule{0pt}{2.5ex}
    1 \hspace{1pt} & 3 & $\tfrac{1}{3}(2k+1)$ & 0 & 0 \\[3pt]
    2 & 2 & $\tfrac{1}{3}(2k-1)$ & 2 & 2 \\[3pt]
    0 & 1 & $\tfrac{1}{3}(2k)$ & 1 & \ 1.
    \end{tabular}\vspace{-1em}
    \caption*{}
\end{table}

\vspace{-1em}
Then, as in the proof of \Cref{thm:QP},
\begin{align*}
    \LD_P(k) = \LD_{R, (1,0)}(k) = n_i \cdot \blocks_i(k)+ r_i = \begin{cases}
         2k+1 & \text{if } k \equiv 1 \pmod 3,\\
            \frac{4}{3} k + \frac{4}{3} & \text{if } k \equiv 2 \pmod 3,\\
            \frac{2}{3} k + 1 & \text{if } k \equiv 0 \pmod 3.
    \end{cases}
\end{align*}
    
\end{example}

 We conjecture that counting lattice diameter segments or more generally counting $\ell$-slices of $P$ that contain the most lattice points among all $\ell$-dimensional slices of $P$ have a similar behavior under dilation of $P$. More precisely:
 
\begin{conjecture}
   Let $P$ be a lattice $d$-polytope, and $\ell \in [d]$. Define $\LD_P^{\ell}(k)$ as the number of $\ell$-dimensional slices of $P$ that contain the most lattice points among all $\ell$-dimensional slices of $P$. Then there exists a positive integer $q$ such that $\LD_P^{\ell}:\N_{\geq q} \to \N$ agrees with a quasi-polynomial function of degree $d-\ell$.
\end{conjecture}

A proof would require new ideas. For example, in many parts of our proof we rely on the fact that for every diameter direction, lattice diameter segments exist that pass through vertices. As shown in \Cref{ex:LD-3d-interior} this does not hold when $d \geq 3$. Furthermore, the relation between the number of lattice points on an $\ell$-slice and its normalized volume, as described in \Cref{lem:nvol}, is more subtle in higher dimensions.
\section{Diameter directions and the discrete Borsuk partition problem} \label{sec:directions-Borsuk}

We begin with \Cref{thm:diam-dir} which bounds the number of diameter directions of a lattice polytope. Thereafter, we expand our study of lattice diameters to bounded sets $S \subset \Z^d$. Note that if $S = \conv(S) \cap \Z^d$, then $\ldiam(S)=\ldiam(\conv (S))$, where $\conv (S)$ is a lattice polytope, but our results hold for any discrete set of lattice points $S$. In \Cref{Lemma:MaxDegree}, we prove an upper bound on the number of lattice diameter segments that end at a given point in $S$. Using this, we prove \Cref{Thm:Borsuk}, the discrete Borsuk partition problem. A lemma which will be of much use in this section, from \cite{rabi89}, is the following:

\begin{lemma}[Rabinowitz, 1989]\label{lemma:Rabinowitz}
    Let $S \subseteq \Z^d$ be a bounded set and let $m \in \N$ satisfy $|S| > m^d$. Then $\conv (S)$ has at least $m+1$ collinear points. In other words, if $\ldiam(S)<m$, then $|S| \leq  m^d$.
\end{lemma}

% \begin{proof}
%     Suppose, for the sake of contradiction, that $\lvert S \rvert > 2^d$. Each point in $\Z^d$ has a well-defined parity vector in $\{0,1\}^d$, obtained by reducing each coordinate modulo 2. Since there are only $2^d$ such parity vectors, by the pigeonhole principle, there exist two distinct points $x, y \in S$ with the same parity vector. This implies that $x \equiv y \pmod{2}$, so their average $(x+y)/2$ is an integer point. But then the midpoint lies in $\Z^d$ and is distinct from both $x$ and $y$, contradicting the assumption that $\diam_\Z(S) = 1$, which means no lattice point lies strictly between any two points of $S$. Therefore, we must have $\lvert S \rvert \leq 2^d$.
% \end{proof}

\subsection{Diameter directions of a lattice polytope}

Alarcon \cite[Theorem~4.1]{Alarcon_PhD} showed that the number of diameter directions of a lattice polygon $P$ is at most six if $\ldiam(P)=1$ and at most four if $\ldiam(P)>1$; see \Cref{fig:diam-dir} for an example.


\begin{figure}[ht!]
  \centering
  % Figures/3.1_diam-dir.tex
% (Se asume \usepackage{tikz} en el preámbulo)
% Definiciones de colores si no existen (para evitar "Undefined color")
% \makeatletter
% \@ifundefined{color@lilac}{\definecolor{lilac}{RGB}{200,162,200}}{}
% \@ifundefined{color@darklav}{\definecolor{darklav}{RGB}{102,51,153}}{}
% \@ifundefined{color@lightgreen}{\definecolor{lightgreen}{RGB}{102,187,106}}{}
% \@ifundefined{color@stronggreen}{\definecolor{stronggreen}{RGB}{27,94,32}}{}
% \@ifundefined{color@darkyellow}{\definecolor{darkyellow}{RGB}{184,134,11}}{}
% \@ifundefined{color@stoneblue}{\definecolor{stoneblue}{RGB}{100,149,237}}{}
% \makeatother

\begin{minipage}{0.48\textwidth}
  \centering
  \begin{tikzpicture}[thick, scale=0.7]
    % malla
    \foreach \x in {-1,...,3} {
      \foreach \y in {-1,...,3} {
        \node at (\x,\y) [circle, draw=lightgray, fill=lightgray, minimum width=2pt] {};
      }
    }
    % polígono
    \draw[line width=0.3mm, gray, fill=lightgray, opacity=.5] (0,1) -- (1,0) -- (2,2) -- cycle;

    % rectas (recortadas al cuadro [-1,3]x[-1,3])
    \draw[line width=0.4mm, orange] (-1,1) -- (3,1);                 % horizontal
    \draw[line width=0.4mm, lilac!50!darklav] (1,-1) -- (1,3);        % vertical
    \draw[line width=0.4mm, lightgreen] (-1,-1) -- (3,3);             % SW-NE
    \draw[line width=0.4mm, stronggreen] (2,-1) -- (-1,2);            % SE-NW
    \draw[line width=0.4mm, darkyellow] (-1,0.5) -- (3,2.5);          % paralela arriba
    \draw[line width=0.4mm, stoneblue] (0.5,-1) -- (2.5,3);           % paralela abajo

    % puntos negros
    \foreach \p in {(1,0),(0,1),(2,2),(1,1)}{
      \node at \p [circle, fill=black, minimum width=3.5pt] {};
    }
  \end{tikzpicture}
    \hspace{-8em}
    \end{minipage} 
    \hfill
\begin{minipage}{0.48\textwidth}
    \hspace{-8em}
  \centering
  \begin{tikzpicture}[thick, scale=0.7]
    % malla
    \foreach \x in {-1,...,3} {
      \foreach \y in {-1,...,3} {
        \node at (\x,\y) [circle, draw=lightgray, fill=lightgray, minimum width=2.5pt] {};
      }
    }
    % polígono
    \draw[line width=0.3mm, gray, fill=lightgray, opacity=0.5] (0,0) -- (2,0) -- (2,2) -- (0,2) -- cycle;

    % rectas hasta la frontera
    \draw[line width=0.4mm, orange] (-1,0) -- (3,0);
    \draw[line width=0.4mm, orange] (-1,1) -- (3,1);
    \draw[line width=0.4mm, orange] (-1,2) -- (3,2);
    \draw[line width=0.4mm, lilac!50!darklav] (0,-1) -- (0,3);
    \draw[line width=0.4mm, lilac!50!darklav] (1,-1) -- (1,3);
    \draw[line width=0.4mm, lilac!50!darklav] (2,-1) -- (2,3);
    \draw[line width=0.4mm, lightgreen] (-1,-1) -- (3,3);
    \draw[line width=0.4mm, stronggreen] (3,-1) -- (-1,3);

    % puntos negros
    \foreach \p in {(2,0),(0,2),(2,2),(0,0),(1,0),(0,1),(1,1),(2,1),(1,2)}{
      \node at \p [circle, fill=black, minimum width=3.5pt] {};
    }
  \end{tikzpicture}
\end{minipage}

  \caption{Lattice diameter lines colored by their diameter direction.}
  \label{fig:diam-dir}
\end{figure}


We show an analogous statement to the first part, when $\ldiam(P)=1$, in higher dimensions. Our computational experiments suggest that an even smaller bound holds when $\ldiam(P)>1$, but this remains an open problem.

Define the direction $\dir(\bx,\by)$ between two points $\bx,\by \in \R^d$ as $\by-\bx$ or any of its non-zero multiples. For $X \subseteq \R^d$ (usually discrete) define the set of its directions as
\begin{align*}
    \dir(X):= \lb \dir(\bx_1,\bx_2): \bx_1,\bx_2 \in X, \ \bx_1 \neq \bx_2  \rb.
\end{align*}

\begin{lemma}\label{lemma:hyp-dir}
    Let $X,Y \subseteq \R^d$ be sets contained in parallel hyperplanes $H(\ba,\alpha)$ and $ H(\ba,\beta)$, respectively. Suppose that $\dir(X) \cap \dir(Y) = \emptyset$. Then all $\dir(\bx,\by)$ where $\bx \in X$ and $\by \in Y$ are \mbox{pairwise distinct}. 
\end{lemma}

\begin{proof}
    Suppose there exist two distinct pairs $(\bx_1,\by_1), (\bx_2,\by_2) \in X \times Y$ with $\dir(\bx_1,\by_1)=\dir(\bx_2,\by_2)$, that is, $\bx_1-\by_1=\lambda(\bx_2-\by_2)$ for some non-zero $\lambda \in \R$. Then, since $\bx_i \in X \subseteq H(\ba,\alpha)$ and $\by_i \in Y \subseteq H(\ba,\beta)$ we have
    \begin{align*}
        0 \neq \alpha - \beta = \langle a, \bx_1 - \by_1 \rangle = \langle a, \lambda(\bx_2 - \by_2) \rangle 
        = \lambda \langle \ba, \bx_2 - 
        \by_2 \rangle 
        =\lambda (\alpha - \beta)
    \end{align*}
    and thus $\lambda =1$. It follows that $\bx_1-\bx_2 =\by_1-\by_2 \in \dir(X) \cap \dir(Y)$, a contradiction.
\end{proof}

\begin{theorem}\label{thm:diam-dir}
    Let $P$ be a lattice $d$-polytope, with $\ldiam(P)=1$. Then $P$ has at most $\binom{2^d}{2}$ diameter directions and this bound is best possible.
\end{theorem}

\begin{proof}
Let $d \in \N$. For $d=1$ the claim is clear, so assume $d \geq 2$. Since $\ldiam(P)<2$, by \ref{lemma:Rabinowitz} $|P \cap \Z^d| \leq 2^d$. A lattice diameter line is determined by two lattice points in $P$, so the number of possible diameter directions is bounded above by $\binom{2^d}{2}$.

Next, we give a construction of a lattice $d$-polytope attaining this bound. For $t \in \N$ and $x \in \N_{\geq 3}$, let\vspace{-5pt}
    \begin{align*}
        \bv_1=(1,t), \ \bv_2=(x, tx+1), \ \bv_3=(-x-1, -tx-t-1)
    \end{align*}
    and let $T_{t,x} = \conv (\{\bv_1,\bv_2,\bv_3\})$.
    Then $(0,0) = \frac{1}{3}\bv_1+ \frac{1}{3}\bv_2+ \frac{1}{3}\bv_3 \in T_{t,x}$ and since  $T_{t,x}$ has area $\frac{3}{2}$, Pick's Theorem \cite{pick99} implies that $T_{t,x}$ contains no lattice points besides $(0,0)$. Moreover, no three points in $T_{t,x}$ are collinear, thus $\ldiam(T_{t,x})=1$. The slopes of the six diameter lines in $T_{t,x}$ are
    \begin{align*}
    t, \ t + \dfrac{1}{2x + 1}, \ t + \dfrac{1}{x + 2}, \ t + \dfrac{1}{x + 1}, \ t + \dfrac{1}{x}, \text{ and }\ t + \dfrac{1}{x - 1},
    \end{align*}
    which are all distinct and contained in $[t,t+\frac{1}{2}]$.

    %We prove via induction over $k$, for $2 \leq k \leq d$, that there exists a lattice $k$-polytope with $\ldiam(P)=2$ that has $\binom{2^k}{2}$ diameter directions. For $d=k$ this proves the claim.
    \emph{Induction Claim:} For $k \in \{2, \ldots, d\}$ there exist $2^{d-k}$ lattice $k$-polytopes $P_1, \ldots, P_{2^{d-k}}$ such that $|P_i \cap \Z^k| = 2^k$, $P_i$ has $\binom{2^k}{2}$ diameter directions, and $\dir(P_i) \cap \dir(P_j)  = \emptyset$ for all $i,j \in [2^{d-k}], \ i \neq j$.
    
    \emph{Base Case:} For $k=2$ consider $\{ T_{i,3} \st i=1, \ldots, 2^{d-2} \}$. As we showed above $T_{i,3}$ contains $4=2^2$ lattice points, $\ldiam(T_{i,3})=2$ and all lines through two different lattice points have a different diameter direction, thus $T_{i,3}$ has $\binom{4}{2}=6$ diameter directions. Furthermore $\dir(T_{i,3}) \subseteq [i,i+1/2]$, thus the diameter directions of all $T_{i,3}$ for $i =1, \ldots, 2^{d-2}$  are distinct. 

    \emph{Induction Hypothesis:} For $k \in \{2, \ldots, d-1\}$ let $P_1, \ldots, P_{2^{d-k}}$ be lattice $k$-polytopes such that $|P_i \cap \Z^k| = 2^k$, $P_i$ has $\binom{2^k}{2}$ diameter directions, and $\dir(P_i) \cap \dir(P_j)  = \emptyset$ for all $i,j \in [2^{d-k}], \ i \neq j$. In particular, any two pairs of lattice points in $P_i$ define lattice diameter lines \mbox{with distinct slopes.}

    \emph{Induction Step:} Let $P_i$ for $i \in [2^{d-k}]$ be as given in the induction hypothesis. Let $\underline{P_i}:= P_{2i-1} \times \{0\}$ and $\overline{P_{2i}}:= P_{2i} \times \{1\}$ and construct polytopes
    \begin{align*}
        Q_i:= \conv \lp \lb \underline{P_i},\overline{P_{2i}}  \rb \rp\subseteq \R^{k+1}.
    \end{align*}
    Since $\underline{P_i}$ and $\overline{P_{2i}}$ are contained in parallel hyperplanes
    \begin{align*}
        |Q_i \cap \Z^{k+1}| = |\underline{P_i} \cap \Z^{k+1}| +|\overline{P_{2i}} \cap \Z^{k+1}| = 2^{k+1}, 
    \end{align*}
    and $\ldiam(Q_i)=2$. By \Cref{lemma:hyp-dir} all $\dir(\bx,\by)$ for any $\bx \in \underline{P_i}$ and $\by \in \overline{P_{2i}}$ are distinct so any two points in $Q_i$ define different lattice diameter lines with distinct diameter directions, i.e., $Q_i$ has $\binom{2^{k+1}}{2}$ diameter directions. Furthermore, again by \Cref{lemma:hyp-dir}, any two diameter directions in $Q_i$ and $Q_j$ are distinct.
    This proves the induction claim. 
\end{proof}

\begin{example}
    We give an example of our recursive construction using the same idea as in the proof of \mbox{\Cref{thm:diam-dir}}.
    Here $T_0=\conv(\{(0,-3), (1,-4), (5,-5)\})$ and $T_1 = \conv(\{ (3,-3),(4,-2), (2,-7)\})$ each have $6=\binom{2^2}{2}$ distinct diameter directions and thus $P = \conv( T_0 \times \{0\}, T_1 \times \{1\})$ has $\binom{2^3}{2}$ diameter directions. See \Cref{fig:construction-diam-dir}.
\begin{figure}[ht!]
    \centering
    \tikzset{every picture/.style={line width=0.75pt}} %set default line width to 0.75pt        

\begin{tikzpicture}[x=0.9pt, y=0.9pt, yscale=-.7, xscale=.8, transform shape=false]
\begin{scope}
    \clip (0,74) rectangle (350,210);
%uncomment if require: \path (0,299); %set diagram left start at 0, and has height of 299

% % x-axis (behind)
% \draw[color=black] (12,161) -- (355,181);

% % z-axis bottom, dashed inside and top
% \draw[color=black] (190,181) -- (190,208);
% \draw[color=black, dash pattern={on 2.25pt off 2.25pt}] (190,113) -- (190,181);
% \draw[color=black] (190,75) -- (190,113);

% % y-axis
% \draw [color=black] (46,183) -- (331,159);

% x-axis (slope = 1/17), choose x so x-190 is multiple of 17 ⇒ integer y
\draw[color=black] (20,161) -- (360,181);

% y-axis (slope = -1/12), choose x so x-190 is multiple of 12 ⇒ integer y
\draw[color=black] (46,183) -- (346,158);

% z-axis (vertical through the origin)
\draw[color=pink, line width=5pt] (190,60) -- (190,70);
\draw[color=black] (190,75) -- (190,113);
\draw[color=black, dash pattern={on 2.25pt off 2.25pt}] (190,113) -- (190,181);
\draw[color=black] (190,181) -- (190,208);

%BLACK points
\foreach \k in {-10,-8,-6,-4,-2,0,2,4,6,8,10} {
  \fill[black] ({190+12*\k}, {171-\k}) circle (2pt);
}

\foreach \k in {-7,7} {
  \fill[black] ({190+17*\k}, {171+\k}) circle (2pt);
}

% orange top polygon
\draw [color=orange, draw opacity=1, fill=orange, fill opacity=0.34]
  (235,93) -- (176,130) -- (158,115) -- cycle;

% gray polygon rest
\draw [draw opacity=0, fill=lightgray, fill opacity=0.17]
  (235,93) -- (207,185) -- (166,173) -- (158,115) -- cycle;


% Dashed gray lines
\draw[color=gray, dash pattern={on 2.25pt off 2.25pt}] (176,130) -- (204,154);
\draw[color=gray, dash pattern={on 2.25pt off 2.25pt}] (235,93) -- (204,154);

% gray left edges
\draw[color=gray] (158,115) -- (166,173);
\draw[color=gray] (176,130) -- (166,173);
\draw[color=gray] (176,130) -- (207,185);
\draw[color=gray] (235,93) -- (207,185);

% bottom darklav polygon
\draw [draw opacity=0, fill=darklav, fill opacity=0.11]
  (204,154) -- (207,185) -- (166,173) -- cycle;
\draw [color=darklav] (166,173) -- (207,185); %front edge
\draw [color=darklav, dash pattern={on 1.5pt off 1.5pt}] (204,154) -- (207,185); %dashed edges
\draw [color=darklav, dash pattern={on 1.5pt off 1.5pt}] (166,173) -- (204,154);


% DARKLAV VERTICES
\draw[color=darklav, fill=darklav] (190,171) circle (1.6pt); % center
\draw[color=darklav, fill=darklav] (207,185) circle (1.6pt); % front right
\draw[color=darklav, fill=darklav] (204,154) circle (1.6pt); % back right
\draw[color=darklav, fill=darklav] (166,173) circle (1.6pt); % front left



%ORANGE VERTICES
\draw[color=orange, fill=orange] (190,112) circle (1.6pt); % center 
\draw[color=orange, fill=orange] (158,115) circle (1.6pt); % top left
\draw[color=orange, fill=orange] (176,130) circle (1.6pt); % front
\draw[color=orange, fill=orange] (235,93)  circle (1.6pt); % back right

\end{scope}
\end{tikzpicture}


    \caption{The convex hull of two lattice triangles with each $\binom{2^2}{2}$ diameter directions which yields a $3$-dimensional lattice polytope with $\binom{2^3}{2}$ distinct diameter directions.}
    \label{fig:construction-diam-dir}
\end{figure}
%\input{Figures/3.1_diam-dir-triangles}
\end{example}


   




\subsection{The discrete Borsuk partition problem}\label{sec:2.4}

In this section we prove the discrete version of Borsuk's partition problem for lattice diameters. More generally, we study lattice diameters of bounded sets $S \subset \Z^d$. 
% Recall that here the lattice diameter is $\max_{\bx,\by \in S}|[\bx,\by ] \cap \Z^d|-1$. The lattice diameter of $S$ and $\conv S$ need not be the same, see \Cref{fig:Borsuk-sets} for an example:
Notice that the function $f : \Z^d \times \Z^d \to \R$ defined by  
\[
f(\bx, \by) = \nvol([\bx, \by]) = |[\bx, \by] \cap \Z^d| - 1
\]  
    is a semi-metric, i.e., it is positive-definite and symmetric, so $(\Z^d, f)$ is a semi-metric space. Furthermore, $f$ satisfies $f(\lambda \bx, \lambda \by) = |\lambda| \cdot f(\bx,\by)$ for all $\bx,\by \in \Z^d$ and all $\lambda \in \Z$. From this viewpoint, \Cref{Thm:Borsuk} is closely related to the well-known Boltyansky-Gohberg conjecture \cite{boltjansky1985results} which states that any bounded set $S$ in a $d$-dimensional normed space can be partitioned into at most $2^d$ subsets, each having strictly smaller diameter.

Recall that for a bounded set $S$ the \emph{lattice Borsuk number} of $S$, denoted $\beta_\Z(S)$, is the smallest number of subsets into which $S$ can be partitioned, such that each subset has strictly smaller lattice diameter than $\diam_{\Z} (S)$. In \mbox{\Cref{fig:Borsuk-sets}} we show a Borsuk partition of a set $S=P \cap \Z^d$, where $P$ is a lattice polygon, into parts $S_1$ and $S_2$. Necessarily $\ldiam(S_i)< \ldiam(P)$ for $i=1,2$, but 
note that $\ldiam (\conv (S_1)) = \ldiam(P)$.

%Example of Borsuk lattice diameter notion when non convex point set 

\begin{figure}[ht!]
    \centering
        \begin{tikzpicture}[thick, scale=0.8]\vspace{0.5cm}
        %\draw[line width=0.3mm, white](-3,-1.5) -- (6,-1.5);
    
        \foreach \x in {-1,...,5} {
            \foreach \y in {-1,...,1} {
                \node at (\x,\y) [circle, draw=lightgray, fill= lightgray, minimum width=2pt] {};   
            }
        }
    
        % \fill[line width=0.3mm, lightgray, opacity=0.6] (0,-1) -- (0,1) -- (4,0) -- cycle;
        \draw[line width=0.3mm, black] (0,-1) -- (0,1) -- (4,0) -- cycle;

         \foreach \x in {0,...,3} {
                \node at (\x,0) [circle, fill=darklav, minimum width=4pt] {};
            }
        \node at (0,-1) [circle, fill=orange, minimum width=4pt] {};
        \node at (0,1) [circle, fill=orange, minimum width=4pt] {};
        \node at (4,0) [circle, fill=orange, minimum width=4pt] {};

        \draw(2.5,0.75)node[black]{$P$};
        
        \draw(-0.5,0.2)node[darklav]{$S_2$};
        \draw(4.2,-0.5)node[orange]{$S_1$};
        

    \end{tikzpicture} %\vspace{-0.5cm}
    
    \caption{A Borsuk partition $P \cap \Z^d = S_1 \sqcup S_2$ where $\ldiam(S_1) < \ldiam(\conv(S_1))=\ldiam(P)$.}
    \label{fig:Borsuk-sets}
\end{figure}
% \begin{figure}[ht!]
%     \centering
%     % \includegraphics[width=0.4\linewidth]{Figures/2.40_Borsuk.png}
%     %Example of Borsuk lattice diameter notion when non convex point set 

\begin{figure}[ht!]
    \centering
        \begin{tikzpicture}[thick, scale=0.8]\vspace{0.5cm}
        %\draw[line width=0.3mm, white](-3,-1.5) -- (6,-1.5);
    
        \foreach \x in {-1,...,5} {
            \foreach \y in {-1,...,1} {
                \node at (\x,\y) [circle, draw=lightgray, fill= lightgray, minimum width=2pt] {};   
            }
        }
    
        % \fill[line width=0.3mm, lightgray, opacity=0.6] (0,-1) -- (0,1) -- (4,0) -- cycle;
        \draw[line width=0.3mm, black] (0,-1) -- (0,1) -- (4,0) -- cycle;

         \foreach \x in {0,...,3} {
                \node at (\x,0) [circle, fill=darklav, minimum width=4pt] {};
            }
        \node at (0,-1) [circle, fill=orange, minimum width=4pt] {};
        \node at (0,1) [circle, fill=orange, minimum width=4pt] {};
        \node at (4,0) [circle, fill=orange, minimum width=4pt] {};

        \draw(2.5,0.75)node[black]{$P$};
        
        \draw(-0.5,0.2)node[darklav]{$S_2$};
        \draw(4.2,-0.5)node[orange]{$S_1$};
        

    \end{tikzpicture} %\vspace{-0.5cm}
    
    \caption{A Borsuk partition $P \cap \Z^d = S_1 \sqcup S_2$ where $\ldiam(S_1) < \ldiam(\conv(S_1))=\ldiam(P)$.}
    \label{fig:Borsuk-sets}
\end{figure}
%     \caption{A Borsuk partition $P \cap \Z^d = S_1 \sqcup S_2$ where $\ldiam(S_1) < \ldiam(\conv(S_1))=\ldiam(P)$.}
%     \label{fig:Borsuk-sets}
% \end{figure}

Define the \emph{lattice Borsuk graph} of $S$, denoted $\BG_\Z(S)$, as the graph with vertex set $S$ and edges $\{\bx,\by\} \subseteq S$ if and only if $f(\bx,\by)=\diam_\Z(S)$.
%Recall that $\Delta(\BG_\Z(S))$ is the maximum degree of a vertex in $\BG_\Z(S)$.

\begin{remark}\label{Remark:equivalence}
    For a bounded set $S \subset \Z^d$,  $\beta_\Z(S)=\chi(\BG_\Z(S))$, where $\chi$ is the \textit{chromatic number} of $\BG_\Z(S)$, i.e., the minimum number of colors needed to color $S$ so that no two adjacent vertices in $S$ have the same color.
\end{remark}

To find a good bound on the lattice Borsuk number it will be useful to bound $\Delta(\BG_\Z(S))$, the maximum degree of the lattice Borsuk graph. This is precisely what the next lemma does.

\begin{lemma}\label{Lemma:MaxDegree}
    Let $S \subset \Z^d$ be a bounded set and let $\bx \in S$. Then, the number of lattice diameter segments in $S$ that have $\bx$ as an endpoint is at most $2^d - 1$. Therefore, $\Delta(\BG_\Z(S)) \leq 2^d-1$.
\end{lemma}

\begin{proof}
If $|S| <2$, then there is nothing to show, so suppose $|S|\geq 2$.  After an appropriate translation, we may assume that $\bx = \mathbf{0}$. Define
\[
S_{\mathbf{0}} := \{ \bs \in S : f(\mathbf{0}, \bs) = \operatorname{diam}_\Z(S) \}
\]
% as the set of points in $S$ that form a lattice diameter segment with $\mathbf{0}$.
and consider the corresponding set of vectors
\[
S'_{\mathbf{0}} := \left\{ \frac{\bs}{\operatorname{diam}_\Z(S)} : \bs \in S_{\mathbf{0}} \right\}\cup \{\mathbf{0}\}.
\]

Suppose, for contradiction, that $|S_{\mathbf{0}}| > 2^d-1$, then $|S'_{\mathbf{0}}| > 2^d$, and by \ref{lemma:Rabinowitz}, we have $\operatorname{diam}_\Z(S'_{\mathbf{0}}) \geq 2$. Let $\by', \bz' \in S'_{\mathbf{0}}$ such that $f(\by', \bz') \geq 2$. Since all non-zero vectors in $S'_{\mathbf{0}}$ are primitive, neither $\by'$ nor $\bz'$ is $\mathbf{0}$. So for $\by=\ldiam(S)\cdot\by'$ and $\bz=\ldiam(S)\cdot \bz'$ we get
\[
f(\by, \bz) = \ldiam(S) \cdot f(\by', \bz') \geq \ldiam(S) \cdot 2,
\]
which is a contradiction. Thus $|S_\mathbf{0}| \leq 2^d-1$, and $\Delta(\BG_\Z(S)) \leq 2^d-1$ follows from \Cref{Remark:equivalence}.
\end{proof}

Now we state Brooks' Theorem, a graph theoretic result, which relates the maximum degree of a graph to its chromatic number. For the lattice Borsuk graph this translates to a bound on $\beta_\Z(S)$. Brooks' Theorem has many published proofs, see for example \cite[Theorem~8.4]{bondy2008graph}.

\begin{theorem}[Brooks, 1941]\label{Thm:Brooks}
    If $G$ is a connected simple graph, then $\chi(G) \leq \Delta(G)+1$, with equality if and only if $G$ is an odd cycle or a complete graph.
\end{theorem}

Finally, we present a solution to the discrete version of Borsuk's partition problem.

\begin{theorem}\label{Thm:Borsuk}
     Let $S \subset \Z^d$ be a bounded set. Then $\beta_\Z(S) \leq 2^d$ and this bound is best possible.
\end{theorem}

\begin{proof}
    By \Cref{Remark:equivalence}, \Cref{Lemma:MaxDegree} and Brooks' Theorem applied to each component of $\BG_\Z(S)$ we get
    \[
    \beta_\Z(S) = \chi(\BG_\Z(S)) \leq \Delta (\BG_\Z(S)) +1 \leq 2^d.
    \]
    The set $S=\{0,1\}^d$ is an example where this bound is attained.
\end{proof}

Notice that $\beta_\Z(S)= 2^d$ if and only if there exists an $\bx \in S$ that is the endpoint of $2^d-1$ lattice diameter segments, and $\BG_\Z(S)$ is isomorphic to the complete graph on $2^d$ vertices. Thus we conjecture the following:
\begin{conjecture}\label{conj:Borsuk}
    Let $S \subset \Z^d$ be a bounded set. Then $\beta_\Z(S) =2^d$ if and only if $\conv(S)$ is unimodularly equivalent to a $d$-cube $[0,m]^d$ for any $m \in \N$.
\end{conjecture}

In dimension two this conjecture is true and can be proven using a characterization of lattice polygons which have four diameter directions, see \cite[Theorem~2(iii)]{Barany_Furedi}.

\textbf{Acknowledgements}
We are grateful to János Pach, Imre Bárány, Matthias Beck, Ansgar Freyer, Chiara Meroni, and Marie-Charlotte Brandenburg for fruitful discussions, and to Bruce Reznick for providing important references.
This research was partially supported by NSF grants 2348578 and 2434665.

% DO THIS TO COMPILE BIBLIOGRAPHY, for arxiv download the output.bbl file and input it as below (I renamed to refs.bbl)

% \bibliographystyle{amsalpha} { \linespread{1.0} \small \bibliography{references} }

{ \linespread{1.0} \small \input{refs.bbl}}

\vspace{2ex}
% \textsc{Department of Mathematics, University of California, Davis, CA 95616, USA}

\emph{Email addresses:} \texttt{\{aebrose,jadeloera,glopezcampos,antor\}@ucdavis.edu} 


\end{document}